\documentclass[11pt,a4paper]{article}
\usepackage[margin=23mm,headheight=14pt,headsep=8.5mm,footskip=10mm]{geometry}
\usepackage{amsmath,amssymb,amsthm}
\usepackage[no-math]{fontspec}
\usepackage{xeCJK}
\setmainfont{Latin Modern Roman}
\setCJKmainfont{FandolSong-Regular.otf}[Script=CustomDefault,Path=assets/fonts/fandol/,BoldFont=FandolSong-Bold.otf,ItalicFont=FandolSong-Regular.otf,BoldItalicFont=FandolSong-Bold.otf]
\setCJKsansfont{FandolSong-Regular.otf}[Script=CustomDefault,Path=assets/fonts/fandol/,BoldFont=FandolSong-Bold.otf]
\xeCJKsetup{PunctStyle=quanjiao,CheckSingle=true}
\usepackage{setspace}
\setstretch{1.12}
\setlength{\parindent}{2em}
\setlength{\parskip}{3pt}
\usepackage{titlesec}
\titleformat{\section}{\centering\normalfont\Large\bfseries}{\thesection.}{0.6em}{}
\titleformat{\subsection}{\normalfont\large\bfseries}{}{0pt}{}
\titlespacing*{\section}{0pt}{1.8\baselineskip}{1.0\baselineskip}
\titlespacing*{\subsection}{0pt}{1.1\baselineskip}{0.5\baselineskip}
\usepackage{enumitem}
\setlist[enumerate]{leftmargin=2em,labelsep=0.5em,itemsep=3pt,parsep=0pt,topsep=6pt}
\providecommand{\tightlist}{\setlength{\itemsep}{3pt}\setlength{\parskip}{0pt}}
\usepackage{needspace}
\newenvironment{statementblock}{\par\Needspace{9\baselineskip}\noindent}{\par\nopagebreak[4]}
\newenvironment{proofblock}{\par\Needspace{4\baselineskip}\noindent}{\par}
\usepackage{fancyhdr}
\pagestyle{fancy}
\fancyhf{}
\fancyhead[L]{\small 数学手记}
\fancyhead[R]{\small 复变函数 · 第四章}
\fancyfoot[C]{\thepage}
\renewcommand{\headrulewidth}{0.4pt}
\usepackage[hidelinks,unicode,bookmarksnumbered=true,linktoc=all]{hyperref}
\usepackage{bookmark}
\bookmarksetup{depth=2}
\hypersetup{pdftitle={复变函数讲义：第四章},pdfauthor={Vincent Arqain},pdfsubject={全纯函数的 Taylor 展开及其应用}}
\renewcommand{\contentsname}{目录}
\setcounter{tocdepth}{1}
\setcounter{secnumdepth}{1}
\widowpenalty=10000
\clubpenalty=10000
\displaywidowpenalty=10000
\predisplaypenalty=10000
\interdisplaylinepenalty=10000
\raggedbottom

\usepackage{tabularx,booktabs,array}
\renewcommand{\thesection}{4.\arabic{section}}
\emergencystretch=2em

\usepackage{graphicx}
\usepackage{longtable}
\usepackage{calc}
\usepackage{float}
\floatplacement{figure}{H}
\renewcommand{\figurename}{图}
\makeatletter
\def\maxwidth{\ifdim\Gin@nat@width>\linewidth\linewidth\else\Gin@nat@width\fi}
\def\maxheight{\ifdim\Gin@nat@height>0.72\textheight0.72\textheight\else\Gin@nat@height\fi}
\makeatother
\setkeys{Gin}{width=\maxwidth,height=\maxheight,keepaspectratio}
\providecommand{\pandocbounded}[1]{#1}
\setcounter{tocdepth}{2}
\begin{document}
\begin{center}
{\LARGE\bfseries 复变函数讲义\par}
\vspace{8pt}{\large 第四章\quad 全纯函数的 Taylor 展开及其应用\par}
\vspace{6pt}{\small 依史济怀、刘太顺《复变函数》\quad 教材第 103—140 页\par}
\end{center}\vspace{8pt}
本章含5节讲义及全部 91 道习题原题。定义、定理和性质沿用教材的记号与编号；证明补足必要条件、退化情形与关键推导。由紧集上的一致控制建立 Taylor 理论、零点计数与模估计。此次修订补全原函数级数的半径论证、修正 Abel 扇形闭集，并核正习题 4.2.10 的求和起点。习题不附解答，原题中的疑义另列校勘。
\tableofcontents\clearpage
\section{Weierstrass 定理}
对应教材第 103---110 页。本章的结构是：Cauchy
积分公式给出紧集上的导数控制；由此建立级数极限的全纯性，再将积分核展开为几何级数，得到
Taylor 表示；Taylor
的首个非零项决定零点阶数，辐角原理把阶数加总为边界绕数，最后由开映射导出最大模原理与
Schwarz 引理。

沿用教材记号：\(H(D)\) 为域 \(D\)
上的全纯函数全体，\(B(a,r)=\{z:|z-a|<r\}\)，\(\overline E\)
表示闭包。域是非空连通开集；\(\mathbb N=\{1,2,\ldots\}\)。已有的 Cauchy
积分公式（定理 3.4.1）、高阶导数公式（定理 3.4.3---3.4.4）、Cauchy
不等式（定理 3.5.1）与 Morera 定理（定理
3.5.4）作为前置结果引用。本章不使用第五章的 Laurent 展开或残数定理。

\Needspace{5\baselineskip}
\subsection{一致收敛及其判别}\label{ch04-s1-heading-1}

复数项级数 \(\sum_{n=1}^{\infty}z_n\) 收敛，是指部分和
\(S_n=\sum_{j=1}^n z_j\) 收敛。由 \(\mathbb C\)
的完备性，收敛等价于：对每个 \(\varepsilon>0\)，存在 \(N\)，使
\(n>N\)、\(p\ge1\) 时 \(|z_{n+1}+\cdots+z_{n+p}|<\varepsilon\)。取
\(p=1\) 得必要条件 \(z_n\to0\)；若
\(\sum|z_n|<\infty\)，三角不等式使复级数也满足 Cauchy
准则，故绝对收敛蕴含收敛。反向不成立，例如 \(\sum(-1)^{n-1}/n\)。

\textbf{定义 4.1.1。} 设
\(f_n:E\to\mathbb C\)，\(S_n=\sum_{j=1}^n f_j\)。若对每个
\(\varepsilon>0\)，存在与 \(z\in E\) 无关的 \(N\)，使 \(n>N\) 时
\(|S_n(z)-f(z)|<\varepsilon\) 对全部 \(z\in E\) 成立，就称级数在 \(E\)
上一致收敛于 \(f\)，记作
\(S_n\rightrightarrows f\)。等价地，\(\sup_E|S_n-f|\to0\)。

\textbf{定理 4.1.2（一致 Cauchy 准则）。} 级数在 \(E\)
上一致收敛，当且仅当对每个 \(\varepsilon>0\)，存在 \(N\)，使

\[
\left|\sum_{j=n+1}^{n+p}f_j(z)\right|<\varepsilon
\quad(n>N,\ p\ge1,\ z\in E).
\qquad\text{(4.1.3)}
\]

证明：一致收敛时，以 \(\varepsilon/2\) 控制 \(S_{n+p}-f\) 和 \(S_n-f\)
即得。反之，固定 \(z\)，由完备性得 \(S_n(z)\to f(z)\)；先把尾和控制为
\(\varepsilon/2\)，再令 \(p\to\infty\)，得
\(|f(z)-S_n(z)|\le\varepsilon/2<\varepsilon\)，且 \(N\) 与 \(z\)
无关。这里取极限后只能先得到非严格不等号。

\textbf{定理 4.1.3（Weierstrass 判别法）。} 若 \(|f_n(z)|\le a_n\)
对所有 \(z\in E\) 成立，且 \(\sum a_n<\infty\)，则 \(\sum f_n\) 在 \(E\)
上绝对且一致收敛。

证明：尾和的模不超过 \(\sum_{j=n+1}^{n+p}a_j\)，应用定理 4.1.2。对
\(\sum|f_n|\) 用同一估计即可。这一判别法是充分条件；不能从一致收敛反推
\(\sum\sup_E|f_n|<\infty\)。例如在非空集合 \(E\) 上取与 \(z\) 无关的
\(f_n(z)=(-1)^{n-1}/n\)：部分和一致收敛于 \(\log2\)，而
\(\sum\sup_E|f_n|=\sum1/n\) 发散。

\Needspace{5\baselineskip}
\subsection{连续性与逐项积分}\label{ch04-s1-heading-2}

\textbf{定理 4.1.4。} 连续函数项级数在 \(E\) 上一致收敛，其和函数在
\(E\) 上连续（按相对拓扑理解）。

证明：固定 \(a\in E\) 和 \(\varepsilon>0\)，取一个部分和 \(S_N\) 使
\(\sup_E|f-S_N|<\varepsilon/3\)。由 \(S_N\) 在 \(a\) 连续，足够接近
\(a\) 的 \(z\in E\) 满足 \(|S_N(z)-S_N(a)|<\varepsilon/3\)。于是

\[
|f(z)-f(a)|\le |f(z)-S_N(z)|+|S_N(z)-S_N(a)|+|S_N(a)-f(a)|<\varepsilon.
\]

\textbf{定理 4.1.5。} 若 \(f_n\) 在可求长曲线 \(\gamma\) 上连续，且
\(\sum f_n\rightrightarrows f\) 于 \(\gamma\)，则可逐项积分：

\[
\int_\gamma f(z)\,dz=\sum_{n=1}^{\infty}\int_\gamma f_n(z)\,dz.
\qquad\text{(4.1.4)}
\]

证明：定理 4.1.4 保证 \(f\) 连续；若 \(L\) 是曲线长度，长大不等式给出
\(|\int_\gamma(f-S_N)\,dz|\le L\sup_\gamma|f-S_N|\to0\)。方向与重数沿用
\(\gamma\) 的参数化；不要求曲线闭合，也不要求被积函数全纯。若曲线随
\(N\) 变动，则还需控制其长度。

\Needspace{5\baselineskip}
\subsection{紧集上的导数控制}\label{ch04-s1-heading-3}

\textbf{定义 4.1.6。} 在域 \(D\) 的每个紧子集上一致收敛，称为在 \(D\)
中内闭一致收敛，也称局部一致收敛。两种说法等价：每点有一个邻域上一致收敛时，任意紧集可用有限个这样的邻域覆盖；反之，在每点取闭小圆盘即可。

\textbf{定义 4.1.7。} \(G\Subset D\) 表示 \(\overline G\) 是包含于 \(D\)
的紧集。注意紧的是 \(\overline G\)，开集 \(G\) 本身通常不紧。

例如 \(z^n\to0\) 在 \(B(0,1)\) 中内闭一致，但在整个圆盘上不一致，因为
\(\sup_{|z|<1}|z^n|=1\)。教材的级数
\(1+\sum_{k=1}^{\infty}(z^k-z^{k-1})\) 的部分和正是
\(z^n\)。因此``内闭一致''允许靠近边界时失去统一控制。

\textbf{引理 4.1.8。} 设 \(K\subset G\Subset D\)，其中 \(K\) 紧、\(G\)
开。对 \(f\in H(D)\) 和 \(k\ge1\)，有

\[
\sup_K|f^{(k)}|\le C\sup_G|f|,
\qquad C=\frac{k!}{\rho^k},\quad \rho=d(K,\partial G)>0.
\qquad\text{(4.1.5)}
\]

证明：紧集 \(K\) 到闭集 \(\partial G\) 的距离为正。对
\(a\in K\)，\(B(a,\rho)\subset G\)。由定理
3.5.1，\(|f^{(k)}(a)|\le k!\rho^{-k}\sup_G|f|\)，再对 \(a\)
取上确界。若使用圆周积分，先取 \(r<\rho\)，最后令
\(r\uparrow\rho\)，不把可能碰到 \(\partial G\) 的圆周误当成内部曲线。

\textbf{定理 4.1.9（Weierstrass）。} 设 \(f_n\in H(D)\)，级数
\(\sum f_n\) 在 \(D\) 中内闭一致收敛于 \(f\)。则
\(f\in H(D)\)，且对每个固定的 \(k\ge1\)，

\[
f^{(k)}(z)=\sum_{n=1}^{\infty}f_n^{(k)}(z),
\]

右端仍在 \(D\) 中内闭一致收敛。

证明：任取闭圆盘
\(\overline{B(a,r)}\subset D\)。部分和在该闭盘上一致收敛，故 \(f\)
连续。盘内任意三角形边界 \(\gamma\) 上，由定理 4.1.5 及 Cauchy
定理，\(\int_\gamma f=\lim_N\int_\gamma S_N=0\)。Morera 定理给出
\(f\in H(B(a,r))\)；\(a\) 任意，故 \(f\in H(D)\)。

再固定紧集 \(K\subset D\)。选取有界开邻域 \(G\) 使
\(K\subset G\Subset D\)，例如对 \(K\) 作足够小的有限圆盘覆盖。对全纯函数
\(S_N-f\) 应用引理 4.1.8，得到

\[
\sup_K|S_N^{(k)}-f^{(k)}|
\le C\sup_G|S_N-f|
\le C\sup_{\overline G}|S_N-f|\longrightarrow0.
\]

最后一个极限来自 \(\overline G\) 的紧性。先证极限全纯，再对差应用 Cauchy
估计，避免预先对未知可微的极限求导。若
\(D=\mathbb C\)，也用有限圆盘邻域选取 \(G\)，无需写没有有限意义的
\(d(K,\partial\mathbb C)\)。

\textbf{函数列形式。} 若 \(F_n\in H(D)\) 且 \(F_n\to F\) 内闭一致，则
\(F\in H(D)\)，并有 \(F_n^{(k)}\to F^{(k)}\) 内闭一致。

证明：把上面的 \(S_N\) 换成
\(F_N\)，连续性、三角形积分和导数估计的论证完全不变。等价地，对望远镜级数
\(F_1+\sum_{n\ge1}(F_{n+1}-F_n)\) 使用定理 4.1.9。

复分析中只控制函数值就能控制所有固定阶导数，根本原因是导数有 Cauchy
积分表示。实函数 \(\sin(nx)/n\to0\) 一致，而其导数 \(\cos(nx)\)
不收敛，说明这里的全纯性不可删除。

\Needspace{5\baselineskip}
\subsection{例 4.1.10：半平面上的级数}\label{ch04-s1-heading-4}

设

\[
\zeta(z)=\sum_{n=1}^{\infty}\frac1{n^z},\qquad n^z=e^{z\log n}.
\]

对紧集 \(K\subset\{\operatorname{Re}z>1\}\)，存在 \(\delta>0\) 使
\(\operatorname{Re}z\ge1+\delta\) 于 \(K\)，故
\(|n^{-z}|\le n^{-1-\delta}\)。由判别法和 Weierstrass 定理，\(\zeta\)
在该半平面全纯，并且

\[
\zeta^{(k)}(z)=\sum_{n=1}^{\infty}\frac{(-\log n)^k}{n^z}.
\]

对每个固定
\(k\)，\((\log n)^k\le C_{k,\delta}n^{\delta/2}\)，还可直接用收敛的
\(p\) 级数核对导数级数在 \(K\) 上的一致收敛。这里没有断言级数在边界
\(\operatorname{Re}z=1\) 收敛，也没有使用其全纯开拓。

\Needspace{5\baselineskip}
\subsection{定理索引}\label{ch04-s1-heading-5}

\hyperref[ch04-s1-heading-1]{定义 4.1.1、定理 4.1.2---4.1.3} ·
\hyperref[ch04-s1-heading-2]{定理 4.1.4---4.1.5} · \hyperref[ch04-s1-heading-3]{定义
4.1.6---4.1.7、引理 4.1.8、定理 4.1.9} · \hyperref[ch04-s1-heading-4]{例 4.1.10}

\section{幂级数}
对应教材第 110---118 页。幂级数形如
\(\sum_{n=0}^{\infty}a_n(z-z_0)^n\)；平移 \(w=z-z_0\)
后只需讨论中心为零的情形。

\Needspace{5\baselineskip}
\subsection{收敛半径}\label{ch04-s2-heading-1}

\textbf{定义 4.2.1。} 若级数在 \(|z|<R\) 收敛、在 \(|z|>R\) 发散，则称
\(R\in[0,+\infty]\) 为收敛半径；\(B(0,R)\) 为收敛圆。\(R=0\)
时只在中心收敛；\(R=+\infty\) 时处处收敛。这个定义不决定 \(|z|=R\)
上的敛散性。

\textbf{定理 4.2.2（Cauchy--Hadamard 公式）。} 每个幂级数都有收敛半径

\[
R=\frac1L,\qquad L=\limsup_{n\to\infty}\sqrt[n]{|a_n|},
\qquad \frac10=+\infty,\quad \frac1{+\infty}=0.
\]

证明：若 \(|z|L<1\)，可取 \(q<1\) 使最终
\(\sqrt[n]{|a_nz^n|}\le q\)，故级数被几何级数控制而绝对收敛。若
\(|z|L>1\)，可取 \(q>1\)，使无穷多个 \(n\) 满足
\(|a_nz^n|>q^n\)；通项不趋于零，故发散。当 \(L=0\)
时前一论证适用于每个有限 \(z\)；当 \(L=+\infty\) 时后一论证适用于每个
\(z\ne0\)；\(z=0\) 始终收敛。有限正 \(L\) 的情形即给出 \(R=1/L\)。

上极限不能擅自改为普通极限。例如 \(a_n=3^n\)（偶数
\(n\)）、\(a_n=1\)（奇数 \(n\)），则 \(\sqrt[n]{|a_n|}\)
不收敛，而收敛半径仍为 \(1/3\)。

\Needspace{5\baselineskip}
\subsection{Abel 定理与圆内全纯性}\label{ch04-s2-heading-2}

\textbf{定理 4.2.3（Abel）。} 若 \(\sum a_nz^n\) 在某个 \(z_0\ne0\)
处收敛，则在 \(|z|<|z_0|\) 内闭绝对一致收敛。

证明：收敛蕴含通项有界，故存在 \(M\) 使 \(|a_nz_0^n|\le M\)。对任意紧集
\(K\subset B(0,|z_0|)\)，选 \(r<|z_0|\) 使 \(|z|\le r\) 于 \(K\)，则

\[
|a_nz^n|\le M\left(\frac r{|z_0|}\right)^n\quad(z\in K).
\]

右端可求和，应用定理 4.1.3 即得。实际上只需 \(\{a_nz_0^n\}\) 有界。

\textbf{定理 4.2.4。} 收敛半径 \(R>0\)
的幂级数在收敛圆内表示一个全纯函数。

证明：任给 \(K\subset B(0,R)\) 紧，选 \(r<\rho<R\) 使
\(K\subset\overline{B(0,r)}\)。级数在 \(z=\rho\) 收敛，由定理 4.2.3 在
\(K\) 上一致收敛。各项为多项式，再用定理 4.1.9。

\textbf{逐项求导性质。} 若
\(f(z)=\sum_{n=0}^{\infty}a_n(z-z_0)^n\)，则在 \(|z-z_0|<R\)，

\[
f^{(k)}(z)=\sum_{n=k}^{\infty}\frac{n!}{(n-k)!}a_n(z-z_0)^{n-k},\qquad k\ge1,
\]

各式均内闭一致收敛；导数级数的收敛半径仍为 \(R\)。

证明：求导及一致收敛由 Weierstrass 定理。其半径至少为
\(R\)。若一次导数级数在更大的半径 \(R'>R\) 内收敛，则对任意
\(r<R'\)，\(\sum_{n\ge1}|na_n|r^{n-1}<\infty\)，从而
\(\sum_{n\ge1}|a_n|r^n<\infty\)，与 \(R\)
的定义矛盾。逐次应用得所有固定阶的结论。\(R=+\infty\)
时结论直接成立。原函数级数 \(C+\sum a_n(z-z_0)^{n+1}/(n+1)\) 在每个
\(|z-z_0|=r<R\) 上绝对收敛，因为其绝对值级数不超过
\(|C|+r\sum|a_n|r^n\)。若其半径大于
\(R\)，对它应用已证的求导半径不变结论，就得到原级数也有更大半径，矛盾。因此原函数级数的半径恰为
\(R\)，圆内逐项求导等于 \(f\)，从而可逐项积分。

\Needspace{5\baselineskip}
\subsection{圆周上的三种情形}\label{ch04-s2-heading-3}

\textbf{例 4.2.5。} \(\sum_{n=0}^{\infty}z^n\) 的半径为
\(1\)，单位圆周上处处发散，因为通项的模恒为 \(1\)。

\textbf{例 4.2.6。} \(\sum_{n=1}^{\infty}z^n/n^2\) 的半径为
\(1\)，在闭单位圆盘上绝对一致收敛，因为 \(|z^n/n^2|\le1/n^2\)。

\textbf{例 4.2.7。} \(\sum_{n=1}^{\infty}z^n/n\) 的半径为 \(1\)，在
\(z=1\) 发散，在其他单位圆周点收敛。

证明：半径均由 \(\sqrt[n]{n}\to1\) 得出。对
\(z=e^{i\theta}\ne1\)，有限几何和满足
\(|\sum_{j=m}^n z^j|\le2/|1-z|\)。设
\(A_k=\sum_{j=m}^kz^j\)，分部求和给出

\[
\sum_{k=m}^n\frac{z^k}{k}
=\frac{A_n}{n}+\sum_{k=m}^{n-1}A_k\left(\frac1k-\frac1{k+1}\right),
\qquad
\left|\sum_{k=m}^n\frac{z^k}{k}\right|\le\frac{2}{m|1-z|}.
\]

故由 Cauchy 准则收敛。其绝对值级数是调和级数，所以是条件收敛。教材例
4.2.6---4.2.7 将求和起点印成 \(n=0\)，分母此时无意义；此处明确改为
\(n=1\)。

\Needspace{5\baselineskip}
\subsection{非切向极限与 Abel 第二定理}\label{ch04-s2-heading-4}

\textbf{定义 4.2.8。} 对单位圆周点
\(e^{i\theta_0}\)，在盘内沿以指向圆心的半径为轴、半顶角 \(\alpha<\pi/2\)
的固定角域 \(S_\alpha(e^{i\theta_0})\)
趋近该点，所得极限称为非切向极限。对边界点 \(1\)，可写

\[
S_\alpha(1)=\{z\in B(0,1):|\arg(1-z)|\le\alpha\}.
\]

每个固定 \(\alpha<\pi/2\) 分别要求同一极限；不允许一边趋近边界，一边令
\(\alpha\uparrow\pi/2\)。

\textbf{定理 4.2.9（Abel 第二定理）。} 若
\(f(z)=\sum_{n=0}^{\infty}a_nz^n\) 的收敛半径为 \(1\)，且
\(\sum a_n=S\)，则

\[
\lim_{\substack{z\to1\\z\in S_\alpha(1)}}f(z)=S
\qquad(0\le\alpha<\pi/2).
\qquad\text{(4.2.3)}
\]

证明：记 \(\sigma_{n,p}=\sum_{j=1}^p a_{n+j}\)。给定
\(\varepsilon>0\)，取 \(N\) 使 \(n>N\) 时 \(|\sigma_{n,p}|<\varepsilon\)
对所有 \(p\ge1\) 成立。有限分部求和给出

\[
\begin{aligned}
\sum_{j=1}^p a_{n+j}z^{n+j}
&=(1-z)\sum_{j=1}^{p-1}\sigma_{n,j}z^{n+j}+\sigma_{n,p}z^{n+p},\\
\left|\sum_{j=1}^p a_{n+j}z^{n+j}\right|
&\le\varepsilon\left(\frac{|1-z|}{1-|z|}+1\right),\qquad |z|<1.
\end{aligned}
\qquad\text{(4.2.4)}
\]

要让上界与 \(z\) 无关，正需要非切向条件。写
\(1-z=\rho e^{i\theta}\)，\(|\theta|\le\alpha\)、\(\rho\le\cos\alpha\)，则

\[
\frac{|1-z|}{1-|z|}
=\frac{1+|z|}{2\cos\theta-\rho}
\le\frac{2}{\cos\alpha}.
\]

令
\(E_\alpha=\{1-\rho e^{i\theta}:0\le\rho\le\cos\alpha,\ |\theta|\le\alpha\}\)。这是显式包含端点
\(1\) 的紧集；当 \(\rho>0\)
时，\(|z|^2=1-2\rho\cos\theta+\rho^2<1\)。因此尾和在
\(E_\alpha\setminus\{1\}\) 上有统一小界；在端点 \(z=1\)，尾和就是
\(\sigma_{n,p}\)，同样成立。该闭角形只在 \(1\) 碰到单位圆周：若
\(|z|=1\) 且 \(z\ne1\)，则 \(\rho=2\cos\theta\ge2\cos\alpha\)，与
\(\rho\le\cos\alpha\) 矛盾。由一致 Cauchy
准则，级数在这一集合上一致收敛，和函数在 \(1\) 连续，其值为
\(S\)，结论成立。

关键是用尾部系数的部分和控制振荡，不要求 \(\sum|a_n|\) 收敛。径向极限是
\(\alpha=0\)
的特例；任意切向趋近没有上述距离比的统一上界，不能由本定理直接处理。一般边界点通过旋转平移归结到
\(1\)。

\Needspace{5\baselineskip}
\subsection{例 4.2.10：对数与 Fourier 和式}\label{ch04-s2-heading-5}

令 \(f(z)=\sum_{n=1}^{\infty}z^n/n\)。在 \(|z|<1\) 内逐项求导，得
\(f'(z)=1/(1-z)\)。由于 \(1-z\) 位于右半平面，取在 \(z=0\)
处值为零的对数分支，\((-\log(1-z))'=1/(1-z)\) 且两函数在零点均为零，故

\[
\sum_{n=1}^{\infty}\frac{z^n}{n}=-\log(1-z),\qquad |z|<1.
\]

对 \(0<\theta<2\pi\)，例 4.2.7 保证边界级数收敛，Abel
第二定理把它的和与径向极限连接。因为

\[
1-e^{i\theta}=2\sin\frac\theta2\,e^{i(\theta/2-\pi/2)},
\]

所选分支满足 \(\arg(1-e^{i\theta})=(\theta-\pi)/2\)。取实虚部得到

\[
\begin{aligned}
\sum_{n=1}^{\infty}\frac{\cos n\theta}{n}
&=-\log\left(2\sin\frac\theta2\right),\\
\sum_{n=1}^{\infty}\frac{\sin n\theta}{n}
&=\frac{\pi-\theta}{2}.
\end{aligned}
\]

代入 \(\theta=\pi\) 和 \(\pi/2\)，分别得
\(\sum_{n\ge1}(-1)^{n-1}/n=\log2\)、\(\sum_{k\ge0}(-1)^k/(2k+1)=\pi/4\)。论证顺序是先证明边界级数收敛，再用
Abel
定理认定其和，不能把``径向极限存在''当作``边界级数收敛''的充分条件。

\Needspace{5\baselineskip}
\subsection{定理索引}\label{ch04-s2-heading-6}

\hyperref[ch04-s2-heading-1]{定义 4.2.1、定理 4.2.2} · \hyperref[ch04-s2-heading-2]{定理
4.2.3---4.2.4及逐项求导} · \hyperref[ch04-s2-heading-3]{例 4.2.5---4.2.7} ·
\hyperref[ch04-s2-heading-4]{定义 4.2.8、定理 4.2.9} · \hyperref[ch04-s2-heading-5]{例
4.2.10}

\section{全纯函数的 Taylor 展开}
对应教材第 118---124 页。上一节证明了``收敛幂级数给出全纯函数''；本节用
Cauchy 积分公式证明反向，从而把全纯性与局部幂级数表示等同起来。

\Needspace{5\baselineskip}
\subsection{Taylor 定理与余项估计}\label{ch04-s3-heading-1}

\textbf{定理 4.3.1。} 若 \(f\in H(B(z_0,R))\)，则

\[
f(z)=\sum_{n=0}^{\infty}\frac{f^{(n)}(z_0)}{n!}(z-z_0)^n,
\qquad |z-z_0|<R.
\qquad\text{(4.3.1)}
\]

级数在该圆盘内闭绝对一致收敛，且展开唯一。

证明：先固定 \(0<r<\rho<R\)，沿逆时针圆周
\(\gamma_\rho=\{\zeta:|\zeta-z_0|=\rho\}\) 使用 Cauchy 公式。对
\(|z-z_0|\le r\)，核有展开

\[
\frac1{\zeta-z}
=\sum_{n=0}^{\infty}\frac{(z-z_0)^n}{(\zeta-z_0)^{n+1}}.
\]

记 \(M_\rho=\max_{\gamma_\rho}|f|\)，则第 \(n\) 项乘以 \(f(\zeta)\)
后的模不超过 \((M_\rho/\rho)(r/\rho)^n\)。这是对 \(z\)、\(\zeta\)
同时成立的可求和上界，因此可逐项积分：

\[
\begin{aligned}
f(z)&=\frac1{2\pi i}\int_{\gamma_\rho}\frac{f(\zeta)}{\zeta-z}\,d\zeta\\
&=\sum_{n=0}^{\infty}\left(\frac1{2\pi i}\int_{\gamma_\rho}
\frac{f(\zeta)}{(\zeta-z_0)^{n+1}}\,d\zeta\right)(z-z_0)^n\\
&=\sum_{n=0}^{\infty}\frac{f^{(n)}(z_0)}{n!}(z-z_0)^n.
\end{aligned}
\]

最后一式用第三章的高阶导数公式。\(r<R\)
任意，得到整个圆盘上的等式和内闭绝对一致收敛。若另有
\(f(z)=\sum a_n(z-z_0)^n\)，逐项求 \(k\) 阶导数并代入 \(z_0\)，得
\(f^{(k)}(z_0)=k!a_k\)，所以系数唯一。

\textbf{余项的积分形式。} 由有限几何恒等式可得

\[
\begin{aligned}
f(z)-\sum_{n=0}^N\frac{f^{(n)}(z_0)}{n!}(z-z_0)^n
&=\frac{(z-z_0)^{N+1}}{2\pi i}
\int_{\gamma_\rho}\frac{f(\zeta)}{(\zeta-z_0)^{N+1}(\zeta-z)}\,d\zeta,\\
\sup_{|z-z_0|\le r}|R_N(z)|
&\le\frac{M_\rho}{1-r/\rho}\left(\frac r\rho\right)^{N+1}.
\end{aligned}
\]

证明：把核的有限展开及其余项代入 Cauchy 公式；在圆周上
\(|\zeta-z|\ge\rho-r\)，再用长度 \(2\pi\rho\)
即得。此式同时给出误差衰减与交换求和、积分的依据；并非由``所有阶导数存在''直接推断
Taylor 级数等于原函数。

Taylor 级数的收敛半径至少为已知的
\(R\)，可能更大，例如整函数被限制在一个小圆盘时，仍可有无穷收敛半径。由本定理不能断言在
\(|z-z_0|=R\) 上收敛。

\textbf{定理 4.3.2。} \(f\) 在 \(z_0\) 的某邻域全纯，当且仅当它在
\(z_0\) 的某邻域可表示为收敛幂级数。

证明：正向取闭小圆盘并用定理 4.3.1；反向用定理
4.2.4。这里``在点处全纯''指在该点的邻域全纯，不能替换成仅在该点复可微。比如
\(|z|^2\) 在零点复可微，却不在任何零点邻域全纯。

\Needspace{5\baselineskip}
\subsection{零点阶数与因子分解}\label{ch04-s3-heading-2}

\textbf{定义 4.3.3。} 若 \(f\) 在 \(z_0\) 邻域全纯，且

\[
f(z_0)=f'(z_0)=\cdots=f^{(m-1)}(z_0)=0,
\qquad f^{(m)}(z_0)\ne0,
\]

则称 \(z_0\) 为 \(f\) 的 \(m\) 阶零点。\(m=1\) 时称简单零点。

\textbf{命题 4.3.4。} \(z_0\) 为 \(m\) 阶零点，当且仅当在其邻域内

\[
f(z)=(z-z_0)^m g(z),\qquad g\in H(B(z_0,\varepsilon)),\quad g(z_0)\ne0.
\qquad\text{(4.3.3)}
\]

证明：Taylor 级数从第 \(m\) 项开始，提出 \((z-z_0)^m\)
后剩下的幂级数给出全纯函数 \(g\)，且
\(g(z_0)=f^{(m)}(z_0)/m!\ne0\)。反之，把 \(g\) 在 \(z_0\) 的 Taylor
展开乘以 \((z-z_0)^m\)，首个非零系数在第 \(m\) 项，便得所列导数条件。

因此零点重数来自局部首项；``零点个数''一旦按重数计算，\(m\) 阶零点就贡献
\(m\)。

\Needspace{5\baselineskip}
\subsection{从局部消失到整体唯一性}\label{ch04-s3-heading-3}

\textbf{命题 4.3.5。} 若 \(f\in H(D)\)，并在某个小圆盘上恒为零，则
\(f\equiv0\) 于域 \(D\)。

证明：令 \(A\) 为所有使 \(f\) 在其某个邻域恒为零的点所成集合。\(A\)
非空且开。若 \(a_j\in A\)、\(a_j\to a\in D\)，则对每个
\(k\ge0\)，\(f^{(k)}(a_j)=0\)；由高阶导数连续，得 \(f^{(k)}(a)=0\)。在
\(a\) 的小圆盘作 Taylor 展开即知 \(f\) 恒为零，故 \(a\in A\)。所以 \(A\)
在 \(D\) 中又闭，由连通性 \(A=D\)。这一证明只使用 Taylor
定理与连通性，没有循环引用唯一性定理。

\textbf{命题 4.3.6。} 若 \(f\in H(D)\) 且 \(f\not\equiv0\)，则 \(D\)
中每个零点都是有限阶且孤立的。

证明：若 \(f\) 在零点 \(z_0\) 的所有阶导数皆为零，Taylor 展开使 \(f\)
在一个邻域恒零，再由命题 4.3.5 得全域恒零，矛盾。因此存在最小 \(m\ge1\)
使 \(f^{(m)}(z_0)\ne0\)。由命题 4.3.4，\(f=(z-z_0)^mg\)，而连续性使
\(g\) 在足够小邻域内不为零，故其中只有零点 \(z_0\)。

\textbf{定理 4.3.7（唯一性定理）。} 若
\(f_1,f_2\in H(D)\)，存在两两互异的点 \(z_n\in D\) 满足
\(f_1(z_n)=f_2(z_n)\)，且 \(z_n\to a\in D\)，则 \(f_1\equiv f_2\) 于
\(D\)。

证明：令 \(g=f_1-f_2\)。若 \(g\not\equiv0\)，则其零点孤立；但连续性给出
\(g(a)=0\)，而互异零点 \(z_n\) 在 \(a\)
聚集，矛盾。教材原述未明确互异条件，此处补明；重复取同一个点不能推出恒等。

聚点必须在域内。函数 \(f(z)=\sin(1/(1-z))\) 在单位圆盘全纯，并在
\(z_n=1-1/(n\pi)\) 为零，但不恒零；这里聚点 \(1\)
在边界。连通性同样必要：不连通开集上的结论只能在相应连通分支成立。

\textbf{紧集上的零点有限性。} 若 \(f\not\equiv0\) 于 \(D\)，则任意紧集
\(K\subset D\) 中只有有限个零点。

证明：否则在 \(K\) 中取互异零点列，紧性给出一个仍在 \(K\subset D\)
的聚点，与命题 4.3.6
矛盾。下一节应用到曲线及其内部时，必须确保整个闭域紧含于全纯域。

\Needspace{5\baselineskip}
\subsection{常用 Taylor 展开与分支}\label{ch04-s3-heading-4}

指数、正弦、余弦为整函数。将它们在零点的导数代入定理 4.3.1，得到

\[
\begin{aligned}
e^z&=\sum_{n=0}^{\infty}\frac{z^n}{n!},\\
\cos z&=\sum_{n=0}^{\infty}(-1)^n\frac{z^{2n}}{(2n)!},\\
\sin z&=\sum_{n=0}^{\infty}(-1)^n\frac{z^{2n+1}}{(2n+1)!},\qquad z\in\mathbb C.
\end{aligned}
\]

第一式沿用教材公式（4.3.4）。也可先由收敛半径公式确认级数为整函数，再将实轴上的已知展开通过唯一性定理推广到复平面。

在 \(|z|<1\)，取 \(\log(1+z)\) 在零点值为零的主支。由上一节对数和式代入
\(-z\)，有

\[
\log(1+z)=\sum_{n=1}^{\infty}(-1)^{n-1}\frac{z^n}{n}.
\]

对 \(\alpha\in\mathbb C\)，定义
\((1+z)^\alpha=e^{\alpha\log(1+z)}\)。由链式法则归纳，\(f^{(n)}(0)=\alpha(\alpha-1)\cdots(\alpha-n+1)\)，故

\[
(1+z)^\alpha=\sum_{n=0}^{\infty}\binom\alpha n z^n,
\quad
\binom\alpha0=1,\quad
\binom\alpha n=\frac{\alpha(\alpha-1)\cdots(\alpha-n+1)}{n!}.
\]

若 \(\alpha\) 不是非负整数，系数相邻比的模趋于 \(1\)，半径为 \(1\)；若
\(\alpha\) 是非负整数，级数终止为多项式，半径为无穷。展开式在 \(|z|<1\)
恒成立，但边界敛散性仍须另查。对数及非整数幂的分支必须固定，不能在求导过程中变换。

\Needspace{5\baselineskip}
\subsection{定理索引}\label{ch04-s3-heading-5}

\hyperref[ch04-s3-heading-1]{定理 4.3.1---4.3.2、Taylor 余项} ·
\hyperref[ch04-s3-heading-2]{定义 4.3.3、命题 4.3.4} · \hyperref[ch04-s3-heading-3]{命题
4.3.5---4.3.6、定理 4.3.7} · \hyperref[ch04-s3-heading-4]{初等函数展开}

\section{辐角原理和 Rouché 定理}
对应教材第 124---134 页。本节的曲线 \(\gamma\)
均按给定方向和重数积分。说``\(\gamma\) 及其内部位于 \(D\)''是指闭 Jordan
域包含于开集 \(D\)，因其有界闭，实际紧含于 \(D\)。这一条件比仅有
\(\gamma\subset D\) 强；例如环域中的圆周，其内部可能含有不属于 \(D\)
的孔洞。

\Needspace{5\baselineskip}
\subsection{对数导数与零点计数}\label{ch04-s4-heading-1}

\textbf{定理 4.4.1。} 设 \(f\in H(D)\)，\(\gamma\)
为逆时针的可求长简单闭曲线，且 \(\gamma\) 及其内部位于 \(D\)。若 \(f\)
在 \(\gamma\) 上无零点，内部的互异零点为 \(a_1,\ldots,a_n\)，阶数为
\(\alpha_1,\ldots,\alpha_n\)，则

\[
\frac1{2\pi i}\int_\gamma\frac{f'(z)}{f(z)}\,dz
=\sum_{j=1}^n\alpha_j.
\qquad\text{(4.4.1)}
\]

证明：边界非零排除
\(f\equiv0\)；零点孤立性与闭域紧性保证零点只有有限个。取两两不交的闭小圆盘
\(\overline{B(a_j,\varepsilon_j)}\)，全在 \(\gamma\)
内部，并且在每个盘内 \(f(z)=(z-a_j)^{\alpha_j}g_j(z)\)，\(g_j\)
全纯且无零点。挖去这些圆盘后，\(f'/f\) 全纯。由多连通 Cauchy 定理，

\[
\int_\gamma\frac{f'}f\,dz
=\sum_{j=1}^n\int_{\gamma_j}\frac{f'}f\,dz,
\qquad \gamma_j=\partial B(a_j,\varepsilon_j)\text{ 取逆时针方向}.
\]

这里挖孔域的内边界本来是顺时针，移到等式右边才变为逆时针。又

\[
\frac{f'}f=\frac{\alpha_j}{z-a_j}+\frac{g_j'}{g_j}.
\]

第二项在小圆盘全纯，积分为零；第一项积分为
\(2\pi i\alpha_j\)，求和即得。无零点时直接用 Cauchy
定理得零。整个证明只用局部分解与 Cauchy 定理，不需要残数理论，也不要求
\(f\) 在整个域上有单值对数。

\Needspace{5\baselineskip}
\subsection{辐角增量的准确含义}\label{ch04-s4-heading-2}

设连续曲线
\(\lambda:[a,b]\to\mathbb C\setminus\{0\}\)。可沿参数选取连续辐角
\(\theta(t)\)，使 \(\lambda(t)=|\lambda(t)|e^{i\theta(t)}\)，并定义
\(\Delta\operatorname{Arg}\lambda=\theta(b)-\theta(a)\)。

\textbf{连续辐角的存在与唯一性。} 其存在性可由有限分段构造：紧参数区间上
\(|\lambda|\)
有正下界，利用一致连续性把区间分为有限段，使每一段像落在不含零的小圆盘内；各盘有局部对数，相邻段的辐角加上适当的
\(2\pi\) 整数倍即可拼接。两个连续辐角之差属于
\(2\pi\mathbb Z\)，又连续，故为常数，因而增量与初始分支无关。

若曲线可求长，用弧长参数（Lipschitz、因而绝对连续）逐段对局部对数使用链式法则，得到

\[
\int_\lambda\frac{dw}{w}
=\log\frac{|\lambda(b)|}{|\lambda(a)|}+i\Delta\operatorname{Arg}\lambda.
\]

对闭曲线，模的对数差为零，首尾辐角相差 \(2\pi\) 的整数倍，故

\[
\operatorname{Ind}(\lambda,0)
=\frac1{2\pi i}\int_\lambda\frac{dw}{w}
=\frac1{2\pi}\Delta\operatorname{Arg}\lambda\in\mathbb Z.
\qquad\text{(4.4.4)}
\]

像曲线可能自交，必须保留参数与重数，不能只看图像所围区域。对\textbf{开放}曲线，即使整个像落在一个有单值辐角的圆盘内，辐角增量也只是末端辐角减去初端辐角，通常不等于零。

\textbf{定理 4.4.2（辐角原理）。} 在定理 4.4.1 的条件下，若 \(N\) 是
\(f\) 在 \(\gamma\) 内部的零点数（计重数），则

\[
N=\frac1{2\pi}\Delta_\gamma\operatorname{Arg}f(z)
=\operatorname{Ind}(f\circ\gamma,0).
\qquad\text{(4.4.6)}
\]

证明：\(f\) 在紧曲线邻域导数有界，故 \(f\circ\gamma\) 可求长；换元得
\(\int_{f\circ\gamma}dw/w=\int_\gamma f'/f\,dz\)，再用式（4.4.4）和定理
4.4.1。例如 \((z^2+1)(z-1)^5\) 在 \(|z|<3\) 有 \(7\)
个零点，故正向圆周的像绕原点 \(7\) 圈。

\Needspace{5\baselineskip}
\subsection{Rouché 定理与扰动稳定性}\label{ch04-s4-heading-3}

\textbf{定理 4.4.3（Rouché）。} 设 \(f,g\in H(D)\)，\(\gamma\)
及其内部紧含于 \(D\)，\(\gamma\) 是可求长简单闭曲线。若

\[
|f(z)-g(z)|<|f(z)|\qquad(z\in\gamma),
\qquad\text{(4.4.7)}
\]

则 \(f,g\) 在 \(\gamma\) 内部的零点数相同（计重数）。

证明：不等式首先保证边界上的 \(f,g\) 均非零。边界曲线 \(g/f\) 落在
\(B(1,1)\)，该圆盘有单值对数，故闭曲线 \(g/f\)
的辐角增量为零。从局部对数相减，得
\(\Delta_\gamma\operatorname{Arg}g-\Delta_\gamma\operatorname{Arg}f=0\)，再用辐角原理。这里只在边界取比值，不声称
\(g/f\) 在内部无极点。

等价的常用形式是 \(|h|<|f|\) 于边界时，\(f\) 与 \(f+h\)
同数。严格小于不可直接改为小于等于，例如 \(f(z)=z\)、\(h(z)=1\)
在单位圆周等模，而 \(f+h\) 的零点恰落在边界。

\textbf{定理 4.4.4。} 设 \(f\in H(D)\)，\(w_0=f(z_0)\)，且 \(z_0\) 是
\(f-w_0\) 的 \(m\) 阶零点。对每个充分小的 \(\rho>0\)，存在
\(\delta>0\)，使任意 \(a\in B(w_0,\delta)\) 的方程 \(f(z)=a\) 在
\(B(z_0,\rho)\) 中恰有 \(m\) 个根（计重数）。

证明：取 \(\rho\) 使闭盘紧含于 \(D\)，且 \(f-w_0\) 在闭盘中只有零点
\(z_0\)。令 \(\delta=\min_{|z-z_0|=\rho}|f(z)-w_0|>0\)。若
\(|a-w_0|<\delta\)，则边界上
\(|(f-a)-(f-w_0)|<\delta\le|f-w_0|\)；Rouché
定理给出结论。最小值处使用的是``\(\le\)''，不把定义最小值后的非严格不等号写成严格大于。

\Needspace{5\baselineskip}
\subsection{开映射、单叶性与全纯逆}\label{ch04-s4-heading-4}

\textbf{推论 4.4.5（补明非常值条件）。} 若 \(f\in H(D)\)
且非常数，\(w_0=f(z_0)\)，则对充分小的 \(\rho>0\)，存在 \(\delta>0\)，使
\(f(B(z_0,\rho))\supset B(w_0,\delta)\)。

证明：唯一性定理保证 \(f-w_0\) 的零点 \(z_0\) 有有限阶，应用定理
4.4.4，每个 \(a\in B(w_0,\delta)\)
至少有一个原像。常函数显然不能满足此结论。

\textbf{定理 4.4.6（开映射定理）。}
非常数全纯函数把域映为域，并把其内开集映为开集。

证明：推论 4.4.5 给出每个像点的圆盘邻域，故 \(f(D)\)
开；连续映射保持连通，故 \(f(D)\) 连通。对 \(D\)
内任意开集，在每点取小圆盘；\(f\)
在小圆盘上不可能为常数，否则由唯一性定理在 \(D\)
恒为常数，所以同样得到局部开映射。

\textbf{定理 4.4.7。} 若 \(f\in H(D)\) 单叶，则 \(f'(z)\ne0\) 于 \(D\)。

证明：若 \(f'(z_0)=0\)，则 \(f-f(z_0)\) 在 \(z_0\) 有 \(m\ge2\)
阶零点。由局部分解求导，\(f'\) 在足够小的闭盘内除 \(z_0\)
外无零点。应用定理 4.4.4，选 \(a\ne f(z_0)\) 且充分接近 \(f(z_0)\)，使
\(f-a\) 在盘内有 \(m\) 个根。它们都不等于
\(z_0\)，故各处导数非零，均为简单根。因此至少有两个互异原像，矛盾。必须选
\(a\ne f(z_0)\)，否则计重数的 \(m\) 个根可能只是同一点。

\textbf{定理 4.4.8。} 若 \(f'(z_0)\ne0\)，则 \(f\) 在 \(z_0\)
的某邻域单叶。

证明：定理 4.4.4 中此时 \(m=1\)，故足够接近 \(f(z_0)\) 的每个值在
\(B(z_0,\rho)\) 中只有一个原像。由连续性缩小到
\(B(z_0,\rho_1)\)，使其像包含于
\(B(f(z_0),\delta)\)；其中若两点同像，就与原像唯一性矛盾。导数处处非零并不蕴含全域单叶，例如
\(e^z\)。

\textbf{定理 4.4.9。} 单叶全纯函数 \(f:D\to G=f(D)\) 的反函数全纯，且

\[
(f^{-1})'(w)=\frac1{f'(z)},\qquad w=f(z).
\]

证明：开映射定理使 \(f\) 是连续开双射，故 \(f^{-1}\) 连续。设
\(w_0=f(z_0)\)，当 \(w\to w_0\) 时 \(z=f^{-1}(w)\to z_0\)，于是

\[
\frac{f^{-1}(w)-f^{-1}(w_0)}{w-w_0}
=\left(\frac{f(z)-f(z_0)}{z-z_0}\right)^{-1}
\longrightarrow\frac1{f'(z_0)}.
\]

分母极限非零来自定理 4.4.7。这样的映射称为双全纯映射。

\Needspace{5\baselineskip}
\subsection{Hurwitz 定理及其推论}\label{ch04-s4-heading-5}

\textbf{定理 4.4.10（Hurwitz）。} 若 \(f_n\in H(D)\) 内闭一致收敛于
\(f\not\equiv0\)，且可求长简单闭曲线 \(\gamma\) 及其内部位于
\(D\)，\(f\) 在 \(\gamma\) 上非零，则充分大的 \(n\) 使 \(f_n\) 和 \(f\)
在 \(\gamma\) 内部有相同零点数。

证明：Weierstrass 定理给出 \(f\in H(D)\)。紧边界上
\(\delta=\min_\gamma|f|>0\)，一致收敛保证最终
\(\sup_\gamma|f_n-f|<\delta\)，由 Rouché
定理得结论。这里比较的是固定边界内的总重数，不意味着各个根的位置完全不变。

\textbf{无零点极限。} 若每个 \(f_n\) 在 \(D\)
上无零点，则极限或者恒零，或者处处非零。

证明：若极限不恒零而在 \(a\)
有零点，取小圆盘隔离该零点并使边界无零点；Hurwitz 定理迫使最终的 \(f_n\)
在盘内有零点，矛盾。

\textbf{定理 4.4.11。} 单叶全纯函数列 \(f_n\) 的内闭一致极限
\(f\)，若非常数，则仍单叶全纯。

证明：全纯性已知。若 \(z_1\ne z_2\) 且 \(f(z_1)=f(z_2)=c\)，则
\(F=f-c\not\equiv0\)。取两个不交的小闭圆盘分别隔离 \(z_1,z_2\)，并使
\(F\) 在圆周非零。\(F_n=f_n-c\to F\) 内闭一致，故 Hurwitz 定理保证最终
\(F_n\) 在两个盘内各有至少一个根。这给出 \(f_n\)
的两个互异同值点，矛盾。原证明中的``各有一个''只需且应写成``各有至少一个''，不能预设极限零点为简单根。常数极限确有可能，例如
\(f_n(z)=z/n\)。

\Needspace{5\baselineskip}
\subsection{例题与边界检查}\label{ch04-s4-heading-6}

\textbf{例 4.4.12。} \(P(z)=z^8-4z^5+z^2-1\) 在单位圆盘内有 \(5\)
个零点。

解：在 \(|z|=1\)，\(|z^8+z^2-1|\le3<4=|-4z^5|\)，故 \(P\) 与 \(-4z^5\)
同数。严格不等式同时排除边界零点。

\textbf{例 4.4.13。} \(P(z)=z^4-6z+3\) 在 \(1<|z|<2\) 内有 \(3\)
个零点。

解：在 \(|z|=1\)，\(|z^4+3|\le4<6=|-6z|\)，盘内一根且圆周无根；在
\(|z|=2\)，\(|-6z+3|\le15<16=|z^4|\)，盘内四根且圆周无根。相减得三根。教材中此处一行把
\(z^4\) 误印为 \(z^2\)，本式按原方程更正。

\textbf{例 4.4.14。} \(P(z)=z^4+2z^3-2z+10\) 在每个象限各有一个根。

证明：先核查坐标轴。对实数 \(x\)，

\[
P(x)=(x^2+x-1)^2+x^2+9>0.
\]

故实轴无根。又 \(P(iy)=y^4+10-2iy(y^2+1)\)，其实部恒正，虚轴也无根。

取第一象限的正向围道：实轴 \(0\to R\)、圆弧
\(Re^{it}\)（\(0\le t\le\pi/2\)）、虚轴 \(iR\to0\)。令

\[
Q(z)=\frac{P(z)}{z^4}=1+\frac2z-\frac2{z^3}+\frac{10}{z^4}.
\]

对充分大的 \(R\)，圆弧上
\(|Q-1|\le2/R+2/R^3+10/R^4<1\)，可取连续主辐角。令
\(\beta_R=\arg Q(iR)=\arg P(iR)\in(-\pi/2,0)\)，而
\(\arg Q(R)=0\)。沿三段的辐角增量依次为

\[
0,\qquad 2\pi+\beta_R,\qquad-\beta_R.
\]

总增量准确等于 \(2\pi\)，故第一象限围道内恰一根。可令 \(R\)
任意增大，所以整个第一象限恰一根。实系数使第四象限也一根；四次多项式共四根且坐标轴无根，剩下两根互为共轭，分别落在第二、第三象限。

教材把圆弧上的 \(\Delta\operatorname{Arg}Q\)
与虚轴段的增量分别写成零；它们实际为相反的非零数，只有相加才抵消。上面保留端点辐角，修正了这一缺口。

\Needspace{5\baselineskip}
\subsection{定理索引}\label{ch04-s4-heading-7}

\hyperref[ch04-s4-heading-1]{定理 4.4.1} · \hyperref[ch04-s4-heading-2]{定理 4.4.2} ·
\hyperref[ch04-s4-heading-3]{定理 4.4.3---4.4.4} · \hyperref[ch04-s4-heading-4]{推论
4.4.5、定理 4.4.6---4.4.9} · \hyperref[ch04-s4-heading-5]{定理 4.4.10---4.4.11}
· \hyperref[ch04-s4-heading-6]{例 4.4.12---4.4.14}

\section{最大模原理和 Schwarz 引理}
对应教材第 134---140
页。前一节把局部非退化性转化为开映射；这一几何性质禁止全纯函数的模在内点出现局部最大值。再把零点因子除去，就得到
Schwarz 引理以及圆盘映射的刚性。

\Needspace{5\baselineskip}
\subsection{最大模原理}\label{ch04-s5-heading-1}

\textbf{定理 4.5.1。} 若 \(f\in H(D)\) 非常数，则 \(|f|\) 不能在 \(D\)
中取得最大值；更强地，它不能在任何内点取得局部最大值。

证明：设 \(|f|\) 在 \(z_0\) 的某邻域不超过
\(|f(z_0)|\)。取其中一个小圆盘 \(U\)。由唯一性定理，\(f|_U\)
仍非常数；开映射定理说明 \(f(U)\) 含有 \(w_0=f(z_0)\)
的小圆盘。这个圆盘一定含模大于 \(|w_0|\) 的点：\(w_0\ne0\)
时沿其径向外移，\(w_0=0\) 时取任一非零点。由原像存在性得到矛盾。

\textbf{定理 4.5.2。} 若 \(D\)
为有界域，\(f\in H(D)\cap C(\overline D)\) 且非常数，则
\(\max_{\overline D}|f|\) 存在，且仅能在边界上达到。因此

\[
|f(z)|\le\max_{\partial D}|f|\quad(z\in\overline D),
\]

对内部点实际严格小于右端。若 \(f\) 为常数，非严格不等式仍成立。

证明：\(\overline D\) 紧且 \(|f|\) 连续，故最大值存在；定理 4.5.1
排除内部极大点。这里不要求边界光滑或 \(D\) 单连通。

\textbf{有界性不能删除。} 在带域 \(D=\{z:|\operatorname{Im}z|<\pi/2\}\)
上，取 \(f(z)=e^{e^z}\)。边界 \(z=x\pm i\pi/2\) 上 \(|f|=1\)，而实轴上
\(|f(x)|=e^{e^x}\to\infty\)。因此仅有边界界不能控制无界域内部。

\textbf{无零点时的最小模性质。} 非常数全纯函数若在 \(D\) 无零点，则
\(|f|\) 不可能在内部取得局部最小值。

证明：\(1/f\) 全纯且非常数；\(|f|\) 的局部正最小值对应 \(|1/f|\)
的局部最大值，与定理 4.5.1 矛盾。允许零点时，最小模零当然可在内部取得。

\Needspace{5\baselineskip}
\subsection{Schwarz 引理与等号刚性}\label{ch04-s5-heading-2}

\textbf{定理 4.5.3（Schwarz）。} 设
\(f\in H(B(0,1))\)，\(|f(z)|\le1\)，且 \(f(0)=0\)，则

\[
|f(z)|\le|z|,\qquad |f'(0)|\le1.
\]

若某个 \(z_0\ne0\) 满足 \(|f(z_0)|=|z_0|\)，或者 \(|f'(0)|=1\)，则存在
\(\theta\in\mathbb R\)，使
\(f(z)=e^{i\theta}z\)；反之，旋转映射均取得等号。

证明：由 Taylor 展开，

\[
g(z)=\begin{cases}f(z)/z,&z\ne0,\\f'(0),&z=0\end{cases}
\]

在整个单位圆盘全纯。任取 \(0<r<1\)，在 \(|z|=r\) 上
\(|g(z)|\le1/r\)，由最大模原理，闭盘内也成立。固定 \(z\) 后取
\(|z|<r<1\) 并令 \(r\uparrow1\)，得 \(|g(z)|\le1\)，从而得到两条不等式。

若任一等号成立，则 \(|g|\) 在某个内点取到最大值 \(1\)，故 \(g\) 为模等于
\(1\) 的常数，写为 \(e^{i\theta}\) 即得。逆向直接代入。整个论证只在
\(r<1\) 的圆周上应用最大模原理，未假定 \(f\) 在单位圆周连续。

\textbf{缩放形式。} 若 \(f:B(0,R)\to B(0,M)\) 全纯且 \(f(0)=0\)，其中
\(R,M>0\)，则 \(|f(z)|\le(M/R)|z|\)、\(|f'(0)|\le M/R\)。

证明：将 Schwarz 引理用于
\(F(\zeta)=f(R\zeta)/M\)；非零点或导数的等号成立当且仅当
\(f(z)=(M/R)e^{i\theta}z\)。这个缩放把估计中每个长度因子的来源明确写出。

\Needspace{5\baselineskip}
\subsection{单位圆盘的全纯自同构}\label{ch04-s5-heading-3}

\textbf{定义 4.5.4。} 若 \(f:D\to D\) 是双全纯映射，称 \(f\) 为 \(D\)
的全纯自同构，其全体记为 \(\operatorname{Aut}(D)\)。

\textbf{群性质。} 复合保持双全纯性且满足结合律，恒等映射为单位元，定理
4.4.9 保证逆映射全纯，故 \(\operatorname{Aut}(D)\)
在复合下构成群。这逐项验证了群公理。

依教材记

\[
\varphi_a(z)=\frac{a-z}{1-\overline a z},\qquad |a|<1,
\qquad \rho_\theta(z)=e^{i\theta}z.
\]

分母在闭单位圆盘无零点，直接展开分子可得

\[
1-|\varphi_a(z)|^2
=\frac{(1-|a|^2)(1-|z|^2)}{|1-\overline a z|^2}.
\]

因而 \(\varphi_a\) 将单位圆盘映入自身；再代入运算得到
\(\varphi_a\circ\varphi_a(z)=z\)，故它是自同构，交换 \(0\) 与
\(a\)。同样显然 \(\rho_\theta\) 为自同构。这是第二章例 2.5.16
的公式，本节重列计算以便检查共轭位置。

\textbf{定理 4.5.5。} 若 \(f\in\operatorname{Aut}(B(0,1))\) 且
\(f(a)=0\)，则存在实数 \(\theta\) 使

\[
f(z)=e^{i\theta}\frac{a-z}{1-\overline a z}.
\]

证明：令 \(g=f\circ\varphi_a\)，则 \(g\) 与 \(g^{-1}\)
都是固定零点的单位圆盘自同构。分别用 Schwarz 引理，得到

\[
|g'(0)|\le1,\qquad |(g^{-1})'(0)|\le1,
\qquad (g^{-1})'(0)g'(0)=1.
\]

故 \(|g'(0)|=1\)，等号刚性给出 \(g(z)=e^{i\theta}z\)。再复合
\(\varphi_a=\varphi_a^{-1}\)
得结论。反之，上述公式是两个已知自同构的复合，故全体自同构已被穷尽。

\Needspace{5\baselineskip}
\subsection{Schwarz--Pick 引理}\label{ch04-s5-heading-4}

\textbf{定理 4.5.6（Schwarz--Pick）。} 设 \(f:B(0,1)\to B(0,1)\)
全纯，\(a\in B(0,1)\)，\(b=f(a)\)。则

\[
\begin{aligned}
\left|\frac{b-f(z)}{1-\overline b f(z)}\right|
&\le\left|\frac{a-z}{1-\overline a z}\right|,\\
|f'(a)|&\le\frac{1-|b|^2}{1-|a|^2}.
\end{aligned}
\]

若第一式在某个 \(z\ne a\) 处取等号，或第二式取等号，则 \(f\)
为圆盘自同构。反之，自同构在两式中处处取等号。

证明：令 \(g=\varphi_b\circ f\circ\varphi_a\)，则 \(g(0)=0\) 且 \(g\)
把单位圆盘映入自身。Schwarz 引理给出 \(|g(\zeta)|\le|\zeta|\)；取
\(\zeta=\varphi_a(z)\)，并用
\(\varphi_a^2=\operatorname{id}\)，即得第一式。直接求导得

\[
\begin{aligned}
\varphi_a'(z)&=-\frac{1-|a|^2}{(1-\overline a z)^2},\\
\varphi_a'(0)&=-(1-|a|^2),\qquad
\varphi_b'(b)=-\frac1{1-|b|^2},\\
|g'(0)|&=|f'(a)|\frac{1-|a|^2}{1-|b|^2}\le1,
\end{aligned}
\]

得到第二式。第一式在 \(z\ne a\) 处取等号时，对应
\(\zeta\ne0\)；第二式取等号时，\(|g'(0)|=1\)。两种情形均由 Schwarz
等号刚性推出 \(g=\rho_\theta\)，所以
\(f=\varphi_b\circ\rho_\theta\circ\varphi_a\) 是自同构。反之，对自同构
\(f\) 与 \(f^{-1}\)
分别应用两条不等式，可得相反方向的不等式，故均为等式。

若记
\(d_*(a,z)=|(a-z)/(1-\overline a z)|\)，第一式表示全纯圆盘自映射不增加这一量，第二式是其局部导数版本。只在一个非平凡位置取等号，就迫使全映射为自同构；这说明最大模原理不只给出上界，还给出达到上界的严格结构。

\Needspace{5\baselineskip}
\subsection{本章的证明方法与使用顺序}\label{ch04-s5-heading-5}

处理极限，先把所需曲线或圆盘放进一个紧邻域，再取一致界；处理展开，先选
\(r<\rho<R\)，使几何比严格小于 \(1\)；处理零点，先用 Taylor
首项确定重数，再检查边界无零点；处理模估计，先核实域有界或改在内部小圆盘上论证。全章反复使用的是同一个思想：在稍大的集合上获得控制，再把结论收回到内部。

零点计数与最大模原理的两条应用线都来自局部 Taylor 理论。Rouché
定理保持边界扰动下的零点总重数，Schwarz
引理则在归一化后保持模估计并刻画等号。第五章将允许孤立奇点与负次幂，当前的全纯版本无需提前引入那些工具。

\Needspace{5\baselineskip}
\subsection{定理索引}\label{ch04-s5-heading-6}

\hyperref[ch04-s5-heading-1]{定理 4.5.1---4.5.2} · \hyperref[ch04-s5-heading-2]{定理
4.5.3} · \hyperref[ch04-s5-heading-3]{定义 4.5.4、定理 4.5.5} ·
\hyperref[ch04-s5-heading-4]{定理 4.5.6}

\clearpage\section*{第四章习题}\phantomsection\addcontentsline{toc}{section}{第四章习题}
以用户上传的史济怀、刘太顺《复变函数》（向禹重排本）为底本，逐页核对印刷第
108---140 页的五组习题，题数依次为 14、12、15、18、32，共 91
题。题面与讲义分离，不附常规解答；原书有误或条件不足的题目仍保留原式，并在各组末尾单列校勘。校勘中的反例仅用于说明原命题不能直接作答。

\(H(\overline D)\) 按教材惯例表示在包含 \(\overline D\)
的某个开邻域全纯；不等同于 \(H(D)\cap C(\overline D)\)。\(B(\infty,R)\)
沿用教材外圆域记号。闭曲线取正向，多连通域的外边界逆时针、内边界顺时针；特殊约定以题面为准。

\hyperref[ch04-practice-heading-1]{习题 4.1} · \hyperref[ch04-practice-heading-2]{习题 4.2} ·
\hyperref[ch04-practice-heading-3]{习题 4.3} · \hyperref[ch04-practice-heading-4]{习题 4.4} ·
\hyperref[ch04-practice-heading-5]{习题 4.5}

\Needspace{5\baselineskip}
\subsection{习题 4.1}\label{ch04-practice-heading-1}

教材第 108---110 页，共 14 题。

\par\Needspace{5\baselineskip}
\noindent\textbf{1.} 证明：复数项级数 \(\sum_{n=1}^{\infty}z_n\) 收敛，当且仅当
\(\sum_{n=1}^{\infty}\operatorname{Re}z_n\) 和
\(\sum_{n=1}^{\infty}\operatorname{Im}z_n\) 同时收敛。

\par\Needspace{5\baselineskip}
\noindent\textbf{2.} 证明：若 \(\{a_n\}\) 和 \(\{b_n\}\) 满足条件

（1）\(\{\sum_{k=1}^n a_k\}\) 有界；

（2）\(\lim_{n\to\infty}b_n=0\)；

（3）\(\sum_{n=1}^{\infty}|b_n-b_{n+1}|<\infty\)，

则级数 \(\sum_{n=1}^{\infty}a_nb_n\) 收敛，并验证这是 Dirichlet 判别法和
Abel 判别法的推广。

\par\Needspace{5\baselineskip}
\noindent\textbf{3.} 设 \(\sum_{n=1}^{\infty}f_n(z)\) 是非空点集 \(E\)
上的函数项级数。证明：\(\sum_{n=1}^{\infty}f_n(z)\) 在 \(E\)
上一致收敛，当且仅当 \(\sum_{n=1}^{\infty}\operatorname{Re}f_n(z)\) 和
\(\sum_{n=1}^{\infty}\operatorname{Im}f_n(z)\) 都在 \(E\) 上一致收敛。

\par\Needspace{5\baselineskip}
\noindent\textbf{4.} 设
\(0<\alpha<\pi/2\)，\(-\alpha\le\arg z_n\le\alpha\)，\(\forall n\in\mathbb N\)。证明：级数
\(\sum z_n\)、\(\sum\operatorname{Re}z_n\) 和 \(\sum|z_n|\)
具有相同的敛散性。

\par\Needspace{5\baselineskip}
\noindent\textbf{5.} 设
\(\operatorname{Re}z_n\ge0\)，\(\forall n\in\mathbb N\)。证明：若
\(\sum_{n=1}^{\infty}z_n\) 和 \(\sum_{n=1}^{\infty}z_n^2\) 都收敛，则
\(\sum_{n=1}^{\infty}|z_n|^2\) 也收敛。

\par\Needspace{5\baselineskip}
\noindent\textbf{6.} 设 \(\sum_{n=1}^{\infty}z_n\) 是复数项级数，且
\(\limsup_{n\to\infty}\sqrt[n]{|z_n|}=q\)。证明：

（1）若 \(q<1\)，则 \(\sum z_n\) 绝对收敛；

（2）若 \(q>1\)，则 \(\sum z_n\) 发散。

\par\Needspace{5\baselineskip}
\noindent\textbf{7.} 求下列函数项级数的收敛点集：

\[
\text{(1)}\quad\sum_{n=1}^{\infty}\frac{\cos nz}{n^2};
\qquad
\text{(2)}\quad\sum_{n=1}^{\infty}\frac{z^n}{1-z^n}.
\]

\par\Needspace{5\baselineskip}
\noindent\textbf{8.} 设
\(z_n\in\mathbb C\setminus\{0\}\)，\(\forall n\in\mathbb N\)，且
\(\limsup_{n\to\infty}|z_{n+1}/z_n|=q\)。证明：

（1）若 \(q<1\)，则 \(\sum z_n\) 绝对收敛；

（2）若 \(q>1\)，则 \(\sum z_n\) 可能收敛也可能发散。

\textbf{9.（Raabe 判别法）} 设
\(z_n\in\mathbb C\setminus\{0\}\)，\(\forall n\in\mathbb N\)，且
\(\lim_{n\to\infty}|z_{n+1}/z_n|=1\)。证明：若

\[
\limsup_{n\to\infty}n\left(\left|\frac{z_{n+1}}{z_n}\right|-1\right)<-1,
\]

则 \(\sum_{n=1}^{\infty}z_n\) 绝对收敛。

\par\Needspace{5\baselineskip}
\noindent\textbf{10.} 判别下列级数的敛散性：

（1）\(\displaystyle\sum_{n=1}^{\infty}\frac{e^{inz}}n\)，\(\operatorname{Im}z>0\)；

（2）

\[
\sum_{n=1}^{\infty}
\frac{\alpha(\alpha+1)\cdots(\alpha+n-1)\beta(\beta+1)\cdots(\beta+n-1)}
{n!\gamma(\gamma+1)\cdots(\gamma+n-1)},
\]

\(\operatorname{Re}(\alpha+\beta-\gamma)<0\)，\(\gamma\ne0,-1,-2,\ldots\)。

\par\Needspace{5\baselineskip}
\noindent\textbf{11.} 证明：\(\sum_{n=1}^{\infty}(-1)^{n-1}/(n-z)\) 在
\(\mathbb C\setminus\mathbb N\) 上内闭一致收敛。

\par\Needspace{5\baselineskip}
\noindent\textbf{12.} 设 \(\sum_{n=1}^{\infty}f_n(z)\) 是域 \(D\)
上的全纯函数项级数。证明：若 \(\sum\operatorname{Re}f_n(z)\) 在 \(D\)
上内闭一致收敛，则 \(\sum\operatorname{Im}f_n(z)\) 或者在 \(D\)
上内闭一致收敛，或者在 \(D\) 上处处发散。

\par\Needspace{5\baselineskip}
\noindent\textbf{13.} 证明：若域 \(D\) 上的全纯函数列 \(\{f_n(z)\}\) 在 \(D\)
上内闭一致收敛于 \(f(z)\)，则 \(f(z)\) 在 \(D\) 上全纯，并且
\(\{f_n^{(k)}(z)\}\) 在 \(D\) 上内闭一致收敛于
\(f^{(k)}(z)\)，\(\forall k\in\mathbb N\)。

\par\Needspace{5\baselineskip}
\noindent\textbf{14.} 若 \(\{\lambda_n\}\) 是严格单调增加，并且以 \(\infty\)
为极限的正数列，则称 \(\sum_{n=1}^{\infty}a_ne^{-\lambda_nz}\) 为
Dirichlet 级数。证明：

（1）若该级数在 \(z_0=x_0+iy_0\) 处收敛，则它在半平面
\(\{z\in\mathbb C:\operatorname{Re}z>x_0\}\) 上内闭一致收敛；

（2）若该级数在 \(z_0=x_0+iy_0\) 处绝对收敛，则它在闭半平面
\(\{z\in\mathbb C:\operatorname{Re}z\ge x_0\}\) 上绝对一致收敛。

\subsubsection{习题 4.1
校勘与约定}\label{ch04-practice-ux4e60ux9898-41-ux6821ux52d8ux4e0eux7ea6ux5b9a}

第 6、8、9 题分别按原书的上极限符号录入；第 9
题首个比值极限是普通极限。第
7（2）题的收敛点必须先在所有项都有定义的点中讨论，不能包含使某个
\(1-z^n=0\) 的点。第 10（2）题若参数使级数终止，按有限和理解。

\Needspace{5\baselineskip}
\subsection{习题 4.2}\label{ch04-practice-heading-2}

教材第 116---117 页，共 12 题。

\par\Needspace{5\baselineskip}
\noindent\textbf{1.} 设 \(\sum_{n=0}^{\infty}a_nz^n\) 和
\(\sum_{n=0}^{\infty}b_nz^n\) 的收敛半径分别为 \(R_1\) 和
\(R_2\)。证明：

（1）\(\sum_{n=0}^{\infty}(a_n\pm b_n)z^n\) 的收敛半径
\(R\ge\min\{R_1,R_2\}\)；

（2）\(\sum_{n=0}^{\infty}a_nb_nz^n\) 的收敛半径 \(R\ge R_1R_2\)；

（3）\(\sum_{n=0}^{\infty}(\sum_{k=0}^na_kb_{n-k})z^n\) 的收敛半径
\(R\ge\min\{R_1,R_2\}\)。

\par\Needspace{5\baselineskip}
\noindent\textbf{2.} 求下列幂级数的收敛半径：

\[
\begin{aligned}
\text{(1)}&\quad\sum_{n=1}^{\infty}z^{n!};
&\text{(2)}&\quad\sum_{n=0}^{\infty}\frac{z^n}{2^{n^2}};\\
\text{(3)}&\quad\sum_{n=0}^{\infty}[3+(-1)^n]^nz^n;
&\text{(4)}&\quad\sum_{n=0}^{\infty}\frac{n^n}{n!}z^n.
\end{aligned}
\]

\par\Needspace{5\baselineskip}
\noindent\textbf{3.} 证明：若 \(\sum_{n=0}^{\infty}a_nz^n\) 在 \(z_0\ne0\)
处绝对收敛，则它在 \(\overline{B(0,|z_0|)}\) 上绝对一致收敛。

\par\Needspace{5\baselineskip}
\noindent\textbf{4.} 设正数列 \(\{a_n\}\) 单调收敛于零。证明：

（1）\(\sum_{n=1}^{\infty}a_nz^n\) 的收敛半径 \(R\ge1\)；

（2）\(\sum_{n=1}^{\infty}a_nz^n\) 在 \(\partial B(0,R)\setminus\{R\}\)
上处处收敛。

\par\Needspace{5\baselineskip}
\noindent\textbf{5.} 证明（Abel 第二定理的又一说法）：若幂级数
\(f(z)=\sum_{n=0}^{\infty}a_n(z-z_0)^n\) 在多角形域 \(G\)
的每个顶点处都收敛，则它必在 \(\overline G\) 上一致收敛。特别地，\(f\)
在 \(\overline G\) 上连续。

\par\Needspace{5\baselineskip}
\noindent\textbf{6.} 证明：\(\sum_{n=1}^{\infty}(-1)^{[\sqrt n]}z^n/n\)
在其收敛圆周 \(\partial B(0,1)\) 上处处收敛，但不绝对收敛。

\par\Needspace{5\baselineskip}
\noindent\textbf{7.} 证明：若 \(f(z)=\sum_{n=0}^{\infty}a_nz^n\) 是 \(B(0,1)\)
上的有界全纯函数，则 \(\sum_{n=0}^{\infty}|a_n|^2<\infty\)。

\par\Needspace{5\baselineskip}
\noindent\textbf{8.} 设 \(\sum_{n=0}^{\infty}a_nz^n\) 的收敛半径 \(R>0\)。证明：

（1）\(\varphi(z)=\sum_{n=0}^{\infty}(a_n/n!)z^n\) 是整函数；

（2）存在正数 \(M\)，使得

\[
|\varphi^{(n)}(z)|\le\frac{Me^{|z|/R}}{R^n},\qquad z\in\mathbb C,\ n\in\mathbb N.
\]

\par\Needspace{5\baselineskip}
\noindent\textbf{9.} 举例说明 Abel 第二定理的逆不成立。

\par\Needspace{5\baselineskip}
\noindent\textbf{10.} 设 \(f(z)=\sum_{n=0}^{\infty}a_nz^n\) 将 \(B(0,R)\)
一一地映为域 \(G\)。证明：\(G\) 的面积为
\(\pi\sum_{n=1}^{\infty}n|a_n^2|R^{2n}\)。

\par\Needspace{5\baselineskip}
\noindent\textbf{11.} 证明：幂级数 \(\sum_{n=0}^{\infty}a_n(z-z_0)^n\) 在域 \(D\)
上一致收敛，当且仅当它在 \(\overline D\) 上一致收敛。

\par\Needspace{5\baselineskip}
\noindent\textbf{12.} 设幂级数 \(f(z)=\sum_{n=0}^{\infty}a_nz^n\) 的收敛半径为
\(1\)，\(z_0\in\partial B(0,1)\)。证明：若
\(\lim_{n\to\infty}na_n=0\)，并且 \(\lim_{r\to1^-}f(rz_0)\) 存在，则
\(\sum_{n=0}^{\infty}a_nz_0^n\) 收敛于 \(\lim_{r\to1^-}f(rz_0)\)。

\subsubsection{习题 4.2
校勘与约定}\label{ch04-practice-ux4e60ux9898-42-ux6821ux52d8ux4e0eux7ea6ux5b9a}

第 1（2）题若出现
\(0\cdot\infty\)，原式无定义；此退化情形不能据此给出非平凡半径下界。第
2（4）题 \(n=0\) 项的 \(0^0\) 可约定为 \(1\)，不影响收敛半径。

\textbf{第 4（2）题原命题不成立。} 取 \(a_n=2^{-n}\)，则 \(R=2\)，在
\(z=-2\) 通项为 \((-1)^n\)，不趋于零。可确认的结论是单位圆周
\(\partial B(0,1)\setminus\{1\}\) 上收敛；若要保留半径 \(R\)，还须对
\(a_nR^n\) 增加相应条件。

第 6 题 \([\sqrt n]\) 表示整数部分。\textbf{第 8（2）题在精确收敛半径
\(R\) 处条件不足。} 取 \(R=1\)、\(a_n=n\)（\(n\ge1\)），则
\(\varphi^{(n)}(0)=n\)，不存在统一 \(M\)。正确的一般估计应对每个
\(0<r<R\) 写为 \(|\varphi^{(n)}(z)|\le M_r e^{|z|/r}/r^n\)；若另知
\(|a_n|R^n\) 有界，才可用 \(R\)。

第 10 题按教材从 \(n=1\)
开始求和，\(|a_n^2|=|a_n|^2\)；面积允许为无穷。第 12 题按圆内趋近理解
\(r\uparrow1\)，原书简写为 \(r\to1\)。

\Needspace{5\baselineskip}
\subsection{习题 4.3}\label{ch04-practice-heading-3}

教材第 122---124 页，共 15 题。

\par\Needspace{5\baselineskip}
\noindent\textbf{1.} 设 \(D\) 是域，\(a\in D\)，函数 \(f\) 在 \(D\setminus\{a\}\)
上全纯。证明：若 \(\lim_{z\to a}(z-a)f(z)=0\)，则 \(f\) 在 \(D\)
上全纯。

\par\Needspace{5\baselineskip}
\noindent\textbf{2.} 将 \(e^{z/(1-z)}\) 在 \(z=0\) 处展开为幂级数。

\par\Needspace{5\baselineskip}
\noindent\textbf{3.} 证明：

（1）\(|e^z-1|\le e^{|z|}-1\le |z|e^{|z|}\)，\(\forall z\in\mathbb C\)；

（2）\((3-e)|z|<|e^z-1|<(e-1)|z|\)，\(0<|z|<1\)。

\par\Needspace{5\baselineskip}
\noindent\textbf{4.} 设
\(f\in H(B(0,R))\cap C(\overline{B(0,R)})\)，\(S_n(z)=\sum_{k=0}^nf^{(k)}(0)z^k/k!\)。证明：

\[
\begin{aligned}
\text{(1)}\quad S_n(z)&=\frac1{2\pi i}\int_{|\zeta|=R}f(\zeta)
\frac{\zeta^{n+1}-z^{n+1}}{(\zeta-z)\zeta^{n+1}}\,d\zeta,\\
\text{(2)}\quad f(z)-S_n(z)&=\frac{z^{n+1}}{2\pi i}\int_{|\zeta|=R}
\frac{f(\zeta)}{\zeta^{n+1}(\zeta-z)}\,d\zeta,
\qquad z\in B(0,R).
\end{aligned}
\]

\par\Needspace{5\baselineskip}
\noindent\textbf{5.} 是否存在 \(f\in H(B(0,1))\)，使得下述条件之一成立：

（1）\(f(1/n)=n/(n+1)\)，\(n=2,3,4,\ldots\)；

（2）\(f(1/(2n))=0\)，\(f(1/(2n-1))=1\)，\(n=1,2,3,\ldots\)；

（3）\(f(1/n)=f(-1/n)=1/n^2\)，\(n=2,3,4,\ldots\)；

（4）\(f(1/n)=f(-1/n)=1/n^3\)，\(n=2,3,4,\ldots\)。

\par\Needspace{5\baselineskip}
\noindent\textbf{6.} 设 \(f(z)=\sum_{n=0}^{\infty}a_nz^n\) 的收敛半径
\(R>0\)，\(0<r<R\)，\(A(r)=\max_{|z|=r}\operatorname{Re}f(z)\)。证明：

（1）\(\displaystyle a_nr^n=\frac1\pi\int_0^{2\pi}[\operatorname{Re}f(re^{i\theta})]e^{-in\theta}\,d\theta\)，\(\forall n\in\mathbb N\)；

（2）\(|a_n|r^n\le2A(r)-2\operatorname{Re}f(0)\)，\(\forall n\in\mathbb N\)。

\par\Needspace{5\baselineskip}
\noindent\textbf{7.} 设 \(f(z)=1+\sum_{n=1}^{\infty}a_nz^n\) 在 \(B(0,1)\)
上全纯，并且
\(\operatorname{Re}f(z)\ge0\)，\(\forall z\in B(0,1)\)。证明：

（1）\(|a_n|\le2\)，\(\forall n\in\mathbb N\)；

（2）\(\displaystyle\frac{1-|z|}{1+|z|}\le\operatorname{Re}f(z)\le|f(z)|\le\frac{1+|z|}{1-|z|}\)，\(\forall z\in B(0,1)\)；

（3）\(|a_1^2-a_2|\le2\)，\(|2a_1a_2-a_1^3-a_3|\le2\)。

\par\Needspace{5\baselineskip}
\noindent\textbf{8.} 证明：所有实变量的三角恒等式在复变量时也成立。

\par\Needspace{5\baselineskip}
\noindent\textbf{9.} 证明：所有实变量的初等函数的幂级数展开式在复变量时也成立。

\par\Needspace{5\baselineskip}
\noindent\textbf{10.} 证明：若函数 \(1/\cos z\) 在 \(z=0\) 处的 Taylor 级数为
\(\sum_{n=0}^{\infty}(-1)^nE_{2n}z^{2n}/(2n)!\)，则 Euler 数 \(E_{2n}\)
满足关系式

\[
E_0=1,\qquad\sum_{k=0}^n\binom{2n}{2k}E_{2k}=0.
\]

\par\Needspace{5\baselineskip}
\noindent\textbf{11.} 证明：若 \(z/(e^z-1)\) 在 \(z=0\) 处的 Taylor 级数为
\(\sum_{n=0}^{\infty}B_nz^n/n!\)，则 Bernoulli 数 \(B_n\) 满足关系式

\[
B_0=1,\qquad\sum_{k=0}^n\binom{n+1}kB_k=0.
\]

特别地，\(B_1=-1/2\)，\(B_{2n+1}=0\)，\(\forall n\in\mathbb N\)。

\par\Needspace{5\baselineskip}
\noindent\textbf{12.} 证明：若 \(1/(1-z-z^2)\) 在 \(z=0\) 处的 Taylor 级数为
\(\sum_{n=0}^{\infty}a_nz^n\)，则 Fibonacci 数 \(a_n\) 满足关系式

\[
a_0=a_1=1,\qquad a_n=a_{n-1}+a_{n-2},\quad n\ge2.
\]

\par\Needspace{5\baselineskip}
\noindent\textbf{13.} 设 \(D\)
是有界域，\(f\in H(D)\cap C(\overline D)\)。证明：若 \(f\) 在
\(\partial D\) 上不取零值，则 \(f\) 在 \(D\) 中只有有限个零点。

\par\Needspace{5\baselineskip}
\noindent\textbf{14.} 设 \(D\) 是域，\(a\in D\)，\(f\in H(D)\)，并且
\(\sum_{n=0}^{\infty}f^{(n)}(a)\) 收敛。证明：

（1）\(f\) 是整函数；

（2）\(\sum_{n=0}^{\infty}f^{(n)}(z)\) 在 \(\mathbb C\) 上内闭一致收敛。

\par\Needspace{5\baselineskip}
\noindent\textbf{15.} 设 \(f(x)\) 是 \((-R,R)\) 上的 \(C^\infty\)
实函数。证明：\(f(x)\) 能在 \((-R,R)\) 上展开为它在 \(x=0\) 处的 Taylor
级数，当且仅当存在 \([0,R)\) 上的正值函数 \(M(r)\)，使得当
\(n\ge0\)、\(|x|<r<R\) 时，成立不等式

\[
|f^{(n)}(x)|\le\frac{M(r)n!}{(r-|x|)^{n+1}}.
\]

由此可见，实变量的函数能展开为 Taylor 级数的条件是多么苛刻。

\subsubsection{习题 4.3
校勘与约定}\label{ch04-practice-ux4e60ux9898-43-ux6821ux52d8ux4e0eux7ea6ux5b9a}

第 1 题应理解为``可在 \(a\) 处补定义使其全纯''，不是任意指定 \(f(a)\)
后都全纯。第 5（2）题在 \(n=1\) 使用了域外的 \(f(1)\)；应从 \(n\ge2\)
考察，或另加边界定义，这不改变其内部聚点所致的矛盾。

第 8---9
题限于相应函数和分支都有定义的连通域；对数及非整数幂须固定与实变量一致的分支，不能跨分支机械套用。第
10---11 题的递推式从 \(n\ge1\) 起；第 11 题先在零点作可去延拓。第 14
题``\(f\) 是整函数''应准确理解为``\(f\)
存在唯一的整函数延拓''，第二问指该延拓的导数。

\Needspace{5\baselineskip}
\subsection{习题 4.4}\label{ch04-practice-heading-4}

教材第 132---134 页，共 18 题。

\par\Needspace{5\baselineskip}
\noindent\textbf{1.} 设 \(D\) 是由有限条可求长简单闭曲线围成的域。证明：若
\(f,g\in H(\overline D)\)，\(f\) 在 \(\partial D\) 上没有零点，\(f\) 在
\(D\) 中的全部彼此不同的零点为
\(z_1,z_2,\ldots,z_n\)，其相应的阶数分别为 \(k_1,k_2,\ldots,k_n\)，则

\[
\frac1{2\pi i}\int_{\partial D}g(z)\frac{f'(z)}{f(z)}\,dz
=\sum_{j=1}^nk_jg(z_j).
\]

（说明：这是 Cauchy 积分公式和辐角原理的推广。）

\par\Needspace{5\baselineskip}
\noindent\textbf{2.} 利用辐角原理证明代数学的基本定理。

\par\Needspace{5\baselineskip}
\noindent\textbf{3.} 设 \(\lambda>1\)。证明：方程 \(z=\lambda-e^{-z}\) 在右半平面
\(\{z\in\mathbb C:\operatorname{Re}z>0\}\) 中恰有一个根，并且是正实根。

\par\Needspace{5\baselineskip}
\noindent\textbf{4.} 设 \(0<a_0<a_1<\cdots<a_n\)。证明：三角多项式

\[
a_0+a_1\cos\theta+a_2\cos2\theta+\cdots+a_n\cos n\theta
\]

在 \((0,2\pi)\) 中有 \(2n\) 个不同的零点。

（提示：首先证明 \(P_n(z)=\sum_{k=0}^na_kz^k\) 在 \(B(0,1)\) 中有 \(n\)
个根。）

\par\Needspace{5\baselineskip}
\noindent\textbf{5.} 利用 Rouché 定理证明代数学的基本定理。

\par\Needspace{5\baselineskip}
\noindent\textbf{6.} 设 \(0<r<1\)。证明：当 \(n\) 充分大时，多项式
\(1+2z+3z^2+\cdots+nz^{n-1}\) 在 \(B(0,r)\) 中没有根。

\par\Needspace{5\baselineskip}
\noindent\textbf{7.} 设 \(r>0\)。证明：当 \(n\) 充分大时，多项式
\(1+z+z^2/2!+\cdots+z^n/n!\) 在 \(B(0,r)\) 中没有根。

\par\Needspace{5\baselineskip}
\noindent\textbf{8.} 设 \(f(z)\) 在 \(\overline{B(0,1)}\) 上全纯，并且 \(f'(z)\)
在 \(\partial B(0,1)\) 上无零点。证明：当 \(n\)
充分大时，\(F_n(z)=n[f(z+1/n)-f(z)]\) 与 \(f'(z)\) 在 \(B(0,1)\)
中的零点个数相等。

\par\Needspace{5\baselineskip}
\noindent\textbf{9.} 设 \(D\) 是域，\(f_n:D\to\mathbb C\setminus\{0\}\)
是全纯映射，\(\forall n\in\mathbb N\)。证明：若 \(\{f_n\}\) 在 \(D\)
上内闭一致收敛于 \(f\)，则或者 \(f(D)=\{0\}\)，或者
\(f(D)\subset\mathbb C\setminus\{0\}\)。

\par\Needspace{5\baselineskip}
\noindent\textbf{10.} 利用上题的结论证明：若域 \(D\) 上的单叶全纯函数列
\(\{f_n\}\) 在 \(D\) 上内闭一致收敛于 \(f\)，则或者 \(f\) 是常数，或者
\(f\) 也是 \(D\) 上的单叶全纯函数。

\par\Needspace{5\baselineskip}
\noindent\textbf{11.} 求下列全纯函数在 \(B(0,1)\) 中的零点个数：

（1）\(z^9-2z^6+z^2-8z-2\)；

（2）\(2z^5-z^3+3z^2-z+8\)；

（3）\(z^7-5z^4+z^2-2\)；

（4）\(e^z-4z^n+1\)。

\par\Needspace{5\baselineskip}
\noindent\textbf{12.} 证明：若
\(f\in H(B(0,1))\cap C(\overline{B(0,1)})\)，\(f(\overline{B(0,1)})\subset B(0,1)\)，则
\(f(z)\) 在 \(B(0,1)\) 中有唯一的不动点。

\par\Needspace{5\baselineskip}
\noindent\textbf{13.} 设
\(a_1,a_2,\ldots,a_n\in B(0,1)\)，\(f(z)=\prod_{k=1}^n(a_k-z)/(1-\overline{a_k}z)\)。证明：

（1）若 \(b\in B(0,1)\)，则 \(f(z)=b\) 在 \(B(0,1)\) 中恰有 \(n\) 个根；

（2）若 \(b\in B(\infty,1)\)，则 \(f(z)=b\) 在 \(B(\infty,1)\) 中恰有
\(n\) 个根。

\par\Needspace{5\baselineskip}
\noindent\textbf{14.} 设 \(f\in H(\overline{B(0,R)})\)；\(f\) 在
\(\partial B(0,R)\) 上没有零点，在 \(B(0,R)\) 中的零点个数为
\(N\)。证明：

\[
\max_{|z|=R}\operatorname{Re}\left(z\frac{f'(z)}{f(z)}\right)\ge N.
\]

\par\Needspace{5\baselineskip}
\noindent\textbf{15.} 设 \(f\) 是域 \(D\) 上非常数的全纯函数。证明：存在在 \(D\)
中无极限点的点列 \(\{z_n\}\)，使得对每个 \(z\in D\setminus\{z_n\}\)，有
\(f'(z)\ne0\)。

\par\Needspace{5\baselineskip}
\noindent\textbf{16.} 设 \(D\)
是由可求长简单闭曲线围成的单连通域，\(f\in H(D)\cap C(\overline D)\)。证明：若
\(f\) 在 \(\partial D\) 上取实值，则 \(f\)
为常值函数。举例说明，对于一般的单连通域 \(D\)，结论不再成立。

\textbf{17.（边界对应原理的特例）} 设 \(D\) 是由可求长简单闭曲线
\(\gamma\) 围成的单连通域，\(f\in H(D)\cap C(\overline D)\)。证明：若
\(f\) 将 \(\gamma\) 一一地映为简单闭曲线 \(\Gamma\)，则 \(f\) 将 \(D\)
双全纯地映为由 \(\Gamma\) 围成的单连通域 \(G\)。

\textbf{18.（辐角原理）} 设 \(D\)
是由有限条可求长简单闭曲线围成的域，\(f\in H(D)\cap C(\overline D)\)。证明：若
\(f\) 在 \(\partial D\) 上不取零值，则 \(f\) 在 \(D\) 中的零点个数为

\[
\frac1{2\pi}\Delta_{\partial D}\operatorname{Arg}f(z),
\]

其中，\(n\) 阶零点视为 \(n\) 个零点。

\subsubsection{习题 4.4
校勘与约定}\label{ch04-practice-ux4e60ux9898-44-ux6821ux52d8ux4e0eux7ea6ux5b9a}

第 1、8、14 题中的闭域全纯均指开邻域全纯，这保证边界导数和第 8
题的平移有定义。第 11（4）题按 \(n\in\mathbb N\) 理解。第 13
题按重数计根；外圆域题在扩充平面中理解，无穷远点可能贡献根或极点，不能先排除它。第
15 题零点集允许为空或有限，故``点列''准确说是至多可数的离散点集。

第 17---18 题只有边界连续，不能直接在边界使用 \(f'/f\)
积分公式；应通过内部逼近或连续曲线绕数处理，特别是像曲线不一定可求长。这是习题比正文邻域全纯版本多出的证明要求。这里的``围成的域''按教材指有界内部域；若允许外部无界域，应另处理无穷远点。

\Needspace{5\baselineskip}
\subsection{习题 4.5}\label{ch04-practice-heading-5}

教材第 137---140 页，共 32 题。

\par\Needspace{5\baselineskip}
\noindent\textbf{1.} 设 \(D\)
是域，\(f_n\in H(D)\cap C(\overline D)\)，\(\forall n\in\mathbb N\)。证明：若
\(\sum_{n=1}^{\infty}f_n(z)\) 在 \(\partial D\) 上一致收敛，则必在 \(D\)
上一致收敛。

\par\Needspace{5\baselineskip}
\noindent\textbf{2.} 设
\(f\in H(B(0,R))\cap C(\overline{B(0,R)})\)，\(M=\max_{|z|=R}|f(z)|\)。证明：若
\(z_0\in B(0,R)\setminus\{0\}\) 是 \(f(z)\) 的零点，则

\[
R|f(0)|\le(M+|f(0)|)|z_0|.
\]

\par\Needspace{5\baselineskip}
\noindent\textbf{3.} 设 \(z_1,z_2,\ldots,z_n\in B(\infty,1)\)。证明：存在
\(z_0\in\partial B(0,1)\)，使得 \(\prod_{k=1}^n|z_0-z_k|>1\)。

\par\Needspace{5\baselineskip}
\noindent\textbf{4.} 设 \(f\in H(B(0,1))\)。证明：\(M(r)=\max_{|z|=r}|f(z)|\) 是
\([0,R)\) 上的增函数。

\par\Needspace{5\baselineskip}
\noindent\textbf{5.} 利用最大模原理证明代数学的基本定理。

\par\Needspace{5\baselineskip}
\noindent\textbf{6.} 设
\(f\in H(B(\infty,R))\cap C(\overline{B(\infty,R)})\)，并且
\(\lim_{z\to\infty}f(z)\) 存在。证明：若 \(f\) 非常数，则
\(M(r)=\max_{|z|=r}|f(z)|\) 是 \([R,\infty)\) 上的严格减函数。

\par\Needspace{5\baselineskip}
\noindent\textbf{7.} 设 \(f\) 是域 \(D\) 上非常数的全纯函数。证明：若 \(f\) 在
\(D\) 中没有零点，则 \(|f(z)|\) 在 \(D\) 内不能取得最小值。

\par\Needspace{5\baselineskip}
\noindent\textbf{8.} 设
\(f\in H(B(0,1))\)，\(f(0)=0\)。证明：\(\sum_{n=1}^{\infty}f(z^n)\) 在
\(B(0,1)\) 上绝对且内闭一致收敛。

\textbf{9.（全纯函数的 Hadamard 三圆定理）} 设
\(0<r_1<r_2<\infty\)，\(D=\{z\in\mathbb C:r_1<|z|<r_2\}\)，\(f\in H(D)\cap C(\overline D)\)，\(M(r)=\max_{|z|=r}|f(z)|\)（\(r_1\le r\le r_2\)）。证明：\(\log M(r)\)
在 \([r_1,r_2]\) 上是 \(\log r\) 的凸函数，即

\[
\log M(r)\le
\frac{\log r_2-\log r}{\log r_2-\log r_1}\log M(r_1)
+\frac{\log r-\log r_1}{\log r_2-\log r_1}\log M(r_2).
\]

\par\Needspace{5\baselineskip}
\noindent\textbf{10.} 设
\(f\in H(B(0,R))\)，\(f(B(0,R))\subset B(0,M)\)，\(f(0)=0\)。证明：

（1）\(|f(z)\le(M/R)|z|\)；\(|f'(0)|\le M/R\)，\(\forall z\in B(0,R)\setminus\{0\}\)；

（2）等号成立当且仅当
\(f(z)=(M/R)e^{i\theta}z\)（\(\theta\in\mathbb R\)）。

\par\Needspace{5\baselineskip}
\noindent\textbf{11.} 设 \(f\in H(B(0,1))\)，\(f(0)=0\)，并且存在 \(A>0\)，使得
\(\operatorname{Re}f(z)\le A\)，\(\forall z\in B(0,1)\)。证明：

\[
|f(z)|\le\frac{2A|z|}{1-|z|},\qquad z\in B(0,1).
\]

\textbf{12.（Carathéodory 不等式）} 设
\(f\in H(B(0,R))\cap C(\overline{B(0,R)})\)，\(M(r)=\max_{|z|=r}|f(z)|\)，\(A(r)=\max_{|z|=r}\operatorname{Re}f(z)\)（\(0\le r\le R\)）。证明：

\[
M(r)\le\frac{2r}{R-r}A(R)+\frac{R+r}{R-r}|f(0)|,\qquad r\in[0,R).
\]

\par\Needspace{5\baselineskip}
\noindent\textbf{13.} 设 \(f\in H(B(0,1))\)，\(f(0)=1\)，并且
\(\operatorname{Re}f(z)\ge0\)，\(\forall z\in B(0,1)\)。利用 Schwarz
引理证明：

（1）\(\displaystyle\frac{1-|z|}{1+|z|}\le\operatorname{Re}f(z)\le|f(z)|\le\frac{1+|z|}{1-|z|}\)，\(\forall z\in B(0,1)\)；

（2）等号在 \(z\) 异于零时成立，当且仅当

\[
f(z)=\frac{1+e^{i\theta}z}{1-e^{i\theta}z}\qquad(\theta\in\mathbb R).
\]

\par\Needspace{5\baselineskip}
\noindent\textbf{14.} 设 \(f\in H(B(0,1))\)。证明：存在 \(z_0\in\partial B(0,1)\)
和收敛于 \(z_0\) 的点列 \(\{z_n\}\)，使得 \(\lim_{n\to\infty}f(z_n)\)
存在。

\par\Needspace{5\baselineskip}
\noindent\textbf{15.}
求出所有满足条件``\(|f(z)=1|\)，\(\forall z\in\partial B(0,1)\)''的整函数。

\par\Needspace{5\baselineskip}
\noindent\textbf{16.} 设 \(P_n(z)\) 是 \(n\)
次多项式，\(P_n^*(z)=z^nP_n(1/\overline z)\)。证明：若 \(P_n(z)\)
的所有零点都在 \(B(\infty,1)\) 中，则
\(P_n(z)+e^{i\theta}P_n^*(z)\)（\(\theta\in\mathbb R\)）的零点都在
\(\partial B(0,1)\) 上。

\par\Needspace{5\baselineskip}
\noindent\textbf{17.} 设
\(f\in H(B(0,1))\)，\(f(B(0,1))\subset B(0,1)\)。证明：若
\(z_1,z_2,\ldots,z_n\) 是 \(f\) 在 \(B(0,1)\)
中的所有彼此不同的零点，其阶数分别为 \(k_1,k_2,\ldots,k_n\)，则

\[
|f(z)|\le\prod_{j=1}^n\left|\frac{z_j-z}{1-\overline{z_j}z}\right|^{k_j},
\qquad z\in B(0,1).
\]

特别地，有 \(|f(0)|\le\prod_{j=1}^n|z_j|^{k_j}\)。

\par\Needspace{5\baselineskip}
\noindent\textbf{18.} 设 \(f\in H(B(0,1))\)，\(f(B(0,1))\subset B(0,1)\)。证明：

\[
\frac{\bigl||f(0)|-|z|\bigr|}{1-|f(0)||z|}
\le|f(z)|\le\frac{|f(0)|+|z|}{1+|f(0)||z|}.
\]

\par\Needspace{5\baselineskip}
\noindent\textbf{19.} 设
\(f\in H(B(0,1))\)，\(f(B(0,1))\subset B(0,M)\)。证明：\(M|f'(0)|\le M^2-|f(0)|^2\)。

\par\Needspace{5\baselineskip}
\noindent\textbf{20.} 设
\(f\in H(B(0,1))\)，\(f(0)=0\)，\(f(B(0,1))\subset B(0,1)\)。证明：若存在
\(z_1,z_2\in B(0,1)\)，使得
\(z_1\ne z_2\)，\(|z_1|=|z_2|\)，\(f(z_1)=f(z_2)\)，则

\[
|f(z_1)|=|f(z_2)|\le|z_1|^2=|z_2|^2.
\]

（提示：考虑
\(\displaystyle\left(\frac{f(z_1)-f(z)}{1-\overline{f(z_1)}f(z)}\right)
\left(\frac{1-\overline{z_1}z}{z_1-z}\right)
\left(\frac{1-\overline{z_2}z}{z_2-z}\right)\)。）

\par\Needspace{5\baselineskip}
\noindent\textbf{21.} 设
\(f\in H(B(0,1))\)，\(f(0)=0\)，\(f(B(0,1))\subset B(0,1)\)。证明：

\[
|z|\frac{\bigl||f'(0)|-|z|\bigr|}{1-|f'(0)||z|}
\le|f(z)|\le|z|\frac{|f'(0)|+|z|}{1+|f'(0)||z|}.
\]

\par\Needspace{5\baselineskip}
\noindent\textbf{22.} 设 \(f\in H(B(0,1))\)，\(f(B(0,1))\subset B(0,1)\)。证明：

\[
\left|f(z)-\frac{f(0)(1-|z|^2)}{1-|f(0)|^2|z|^2}\right|
\le\frac{|z|(1-|f(0)|^2)}{1-|f(0)|^2|z|^2}.
\]

\par\Needspace{5\baselineskip}
\noindent\textbf{23.} 设
\(\varphi\in\operatorname{Aut}(\mathbb C_\infty)\)，并且将
\(\mathbb C_\infty\) 中的圆周 \(\gamma\) 仍映为圆周。证明：\(\varphi\)
是分式线性变换。

\par\Needspace{5\baselineskip}
\noindent\textbf{24.} 求出上半平面
\(\mathbb C^+=\{z\in\mathbb C:\operatorname{Im}z>0\}\) 的全纯自同构群
\(\operatorname{Aut}(\mathbb C^+)\)。

\par\Needspace{5\baselineskip}
\noindent\textbf{25.} 设 \(\varphi\) 将 \(B(0,1)\) 双全纯地映为域 \(D\)，\(\psi\)
将 \(B(0,1)\) 双全纯地映为域 \(G\)。证明：若 \(f:D\to G\)
是全纯映射，\(f(\varphi(0))=\psi(0)\)，则

\[
f[\varphi(B(0,r))]\subset\psi(B(0,r))\qquad(0<r<1).
\]

\par\Needspace{5\baselineskip}
\noindent\textbf{26.} 设 \(f\in H(B(0,1))\)，\(f(B(0,1))\subset B(0,1)\)。证明：

\[
|f(z)-f(0)|\le|z|\frac{1-|f(0)|^2}{1-|f(0)||z|}.
\]

\par\Needspace{5\baselineskip}
\noindent\textbf{27.} 设 \(D\) 是以原点 \(O\) 为中心、以 \(z_1,z_2,z_3,z_4\)
为顶点的正方形域，\(f\in H(D)\cap C(\overline D)\)，\(M\) 是 \(|f(z)|\)
在 \(\overline D\) 上的最大值，\(m\) 是 \(|f(z)|\) 在线段 \([z_1,z_2]\)
上的最大值。证明：

（1）\(|f(0)|\le m^{1/4}M^{3/4}\)；

（2）在闭三角形 \(\triangle Oz_1z_2\) 上也有
\(|f(z)|\le m^{1/4}M^{3/4}\)。

\par\Needspace{5\baselineskip}
\noindent\textbf{28.} 设 \(D\)
是下面所述的域，\(f\in H(D)\cap C(\overline D)\)，并且在 \(\overline D\)
上有界。证明：

（1）若
\(D=\{z\in\mathbb C:0<\operatorname{Im}z<1\}\)，\(\lim_{\operatorname{Re}z\to\infty,\ \operatorname{Im}z=0}f(z)=A\)，则

\[
\lim_{\substack{\operatorname{Re}z\to\infty\\0\le\operatorname{Im}z\le r}}f(z)=A
\qquad(0<r<1);
\]

（2）若
\(D=\{z\in\mathbb C:\theta_1<\arg z<\theta_2\}\)（\(-\pi<\theta_1<\theta_2<\pi\)），\(\lim_{|z|\to\infty,\ \arg z=\theta_1}f(z)=A\)，则

\[
\lim_{\substack{|z|\to\infty\\\theta_1\le\arg z\le\theta}}f(z)=A
\qquad(\theta_1<\theta<\theta_2);
\]

（3）\(D=\{z\in\mathbb C:\theta_1<\arg z<\theta_2\}\)（\(-\pi<\theta_1<\theta_2<\pi\)），\(\lim_{|z|\to0,\ \arg z=\theta_1}f(z)=A\)，则

\[
\lim_{\substack{|z|\to0\\\theta_1\le\arg z\le\theta}}f(z)=A
\qquad(\theta_1<\theta<\theta_2).
\]

\par\Needspace{5\baselineskip}
\noindent\textbf{29.} 设 \(D\) 是有界域，\(f\in H(D)\)，\(z_0\in D\)。证明：若
\(f(z_0)=z_0\)，\(f(D)\subset D\)，\(f'(z_0)=1\)，则 \(f(z)\equiv z\)。

\par\Needspace{5\baselineskip}
\noindent\textbf{30.} 设 \(f\in H(B(0,1))\)，\(f(0)=0\)，并且
\(|\operatorname{Re}f(z)|<1\)，\(\forall z\in B(0,1)\)。证明：

（1）\(\displaystyle|\operatorname{Re}f(z)|\le\frac4\pi\arctan|z|\)，\(\forall z\in B(0,1)\)；

（2）\(\displaystyle|\operatorname{Im}f(z)|\le\frac2\pi\log\frac{1+|z|}{1-|z|}\)，\(\forall z\in B(0,1)\)。

\par\Needspace{5\baselineskip}
\noindent\textbf{31.} 设 \(f\) 在 \(B(0,1)\cup\{1\}\)
上全纯，\(f(0)=0\)，\(f(1)=1\)，\(f(B(0,1))\subset B(0,1)\)。证明：\(f'(1)\ge1\)。

\par\Needspace{5\baselineskip}
\noindent\textbf{32.} 设 \(P\) 是一个 \(k\) 次多项式，在单位圆周上满足
\(|P(e^{i\theta})|\le1\)。证明：对任意单位圆盘外的 \(z\)，有
\(|P(z)|\le|z|^k\)。

\subsubsection{习题 4.5
校勘与约定}\label{ch04-practice-ux4e60ux9898-45-ux6821ux52d8ux4e0eux7ea6ux5b9a}

\textbf{第 1 题缺少有界域条件。} 在上半平面取级数
\(\sum_{n\ge1}2^{-n}e^{-iz}\)。它在实轴上一致收敛，但每个非零尾和在域内无界，因此不在整个域上一致收敛。补上
\(D\) 有界后可由最大模原理成立。

第 4 题假设为单位圆盘，结论却出现未定义的 \(R\)；可将结论改为
\([0,1)\)，或将假设改为
\(H(B(0,R))\)。``增函数''按不减理解，常函数允许等号。第 6
题极限指有限极限；若允许扩充平面中的无穷极限，\(f(z)=z\) 就否定结论。第
9 题先排除 \(f\equiv0\)，或另作 \(\log0=-\infty\) 的约定。

第 10（1）题原书漏印 \(|f(z)|\) 的右竖线；第 15
题原书写成``\(|f(z)=1|\)''，应为``\(|f(z)|=1\)''。以上题面保留原误式，作答时须改正。

第
13（2）题若将``等号''理解为链式不等式中任意一个等号，则原述不成立：\(f\equiv1\)
在 \(\operatorname{Re}f=|f|\)
中恒取等号。应指左端下界或右端上界在某个非零点达到；对固定点取等号还须满足相应的相位条件。

\textbf{第 16 题疑似漏印外层共轭。} 原书写
\(P_n^*(z)=z^nP_n(1/\overline z)\)，一般不是多项式；标准反射多项式应为
\(z^n\overline{P_n(1/\overline z)}\)，并在零点按多项式延拓。题面未擅自替换，修订意图待与另一版本核对。

\textbf{第 18、21 题的绝对值下界不成立。} 取 \(f\equiv0\)，在任意非零
\(z\) 已与两题矛盾。第 18 题可确认的正确下界为
\(\max\{0,(|f(0)|-|z|)/(1-|f(0)||z|)\}\)；第 21 题为
\(|z|\max\{0,(|f'(0)|-|z|)/(1-|f'(0)||z|)\}\)。原上界不受此笔误影响。

第 3 题中的 \(z_k\) 须为有限复数，才能使用差 \(z_0-z_k\)。第 17
题既然列出``所有''零点，即另含零点总数有限的假设；有限子集版本也可成立。第
23 题 \(\mathbb C_\infty\) 指扩充平面，圆周按广义圆理解。第 27
题顶点按边界顺序编号，使 \([z_1,z_2]\) 为一条边。第 31 题在 \(1\)
处全纯指存在该点的开邻域全纯延拓；结论同时断言 \(f'(1)\)
为实数。除明确列出的校勘外，未将待核题改写成新的自拟题。

\clearpage\section*{校审记录}\phantomsection\addcontentsline{toc}{section}{校审记录}
校审日期：2026 年 11 月 19
日。下列为原稿校审记录；2026 年 11 月 20 日完成网站稿交接复核，原有数学正文、题面与校勘保留。正式发布记录见网站与仓库。

\subsection{底本与范围}\label{ux5e95ux672cux4e0eux8303ux56f4}

底本为用户上传的史济怀、刘太顺《复变函数》（向禹重排本），印刷第
103---140 页（PDF 第 110---147
页）。逐页读取正文、公式及习题，对闭包横线、复共轭、上极限和边界条件作图像复核。五节依次为
Weierstrass 定理、幂级数、Taylor 展开、辐角原理和 Rouché
定理、最大模原理和 Schwarz 引理。

讲义覆盖教材 47 个编号条目：§4.1 的 4.1.1---4.1.10、§4.2 的
4.2.1---4.2.10、§4.3 的 4.3.1---4.3.7、§4.4 的 4.4.1---4.4.14、§4.5 的
4.5.1---4.5.6，包括定义、引理、命题、定理、推论与例题。非定义类结论均给出证明或完整计算；补充性质另用名称标出，不占用教材编号。

独立习题页收录全部 91 题，分节为 14、12、15、18、32
题，保留小问和教材提示，不附常规解答。每节末尾将原书笔误、条件不足与术语解释单列，避免把更正混成原题。

\subsection{前置依赖与证明检查}\label{ux524dux7f6eux4f9dux8d56ux4e0eux8bc1ux660eux68c0ux67e5}

最新课程目录和文件列表已核对：前三章共有 19
节正文，第四章当时仅有目录。使用第三章已完成的 Cauchy
积分公式、高阶导数公式、Cauchy 估计、Morera 定理和多连通 Cauchy
定理，不重复编写第三章，也未提前调用第五章残数定理。

§4.1 先由局部 Morera
定理证明极限全纯，再估计导数之差。紧性使用闭邻域而非开邻域；整个平面的情形通过有限圆盘覆盖处理。§4.2
区分圆内收敛和边界收敛，用分部求和及距离比证明完整非切向 Abel 定理。§4.3
以两半径控制核展开及余项，先证局部恒零传播，再证零点孤立性与唯一性，避免循环论证。

§4.4
零点有限性由紧性与孤立性共同保证；外、内边界方向明确，像曲线保留重数。Rouché
不等式只在边界取比值。单叶函数导数非零的证明明确选择不同于中心像值的扰动，保证所得多根互异。反函数先证连续，再取差商。Hurwitz
定理固定无零点边界后比较总重数。

§4.5 由开映射证明最大模原理，区分局部最大模与有界域边界最大模。Schwarz
引理只在半径小于 1 的闭盘上取极值；Schwarz--Pick
的共轭、自逆恒等式、链式求导及等号条件逐项核对。

\subsection{已确认的正文校勘}\label{ux5df2ux786eux8ba4ux7684ux6b63ux6587ux6821ux52d8}

\begin{itemize}
\tightlist
\item
  定理 4.1.2 的严格不等式取极限时先得到非严格不等式，采用半个误差修正。
\item
  例 4.2.6---4.2.7 的分母含 \(n\)，求和应从 \(n=1\) 开始；例 4.2.10
  的边界辐角取 \((\theta-\pi)/2\)。
\item
  定理 4.3.7 补明互异点列；聚点须在域内。推论 4.4.5 补明非常数条件。
\item
  定理 4.4.4、4.4.10 的最小值估计保持正确的非严格不等号。定理 4.4.7
  排除扰动目标等于中心值；定理 4.4.11 改用``各有至少一个根''。
\item
  例 4.4.13 中 \(z^4\) 的误印更正。例 4.4.14
  保留开放圆弧的端点辐角，得到 \(2\pi+\beta_R\) 与 \(-\beta_R\)
  的准确抵消，并补证实轴、虚轴均无根。
\end{itemize}

\subsection{习题疑义及待核事项}\label{ux4e60ux9898ux7591ux4e49ux53caux5f85ux6838ux4e8bux9879}

已确认不能直接成立的题目包括：4.2.4（一般收敛圆周结论）、4.2.8（精确半径上的统一系数界）、4.5.1（缺有界域条件）、4.5.18
与 4.5.21（错误绝对值下界）。对应位置给出简短反例及可确认的修正。4.5.13
对等号的笼统表述另作限定；4.3.5（2）、4.5.4 的定义域或符号不一致已注明。

4.5.16
的反射多项式定义疑似漏印外层共轭，原式已忠实保留，并写明标准定义；由于未与另一版教材核对，修订意图保留``待核''，不将推测记为已确认原文。4.5.10
和 4.5.15
的竖线笔误也保留原式并提示更正。其余需补充的退化情形、分支、有限值极限和无穷远点约定均列在各组校勘中。

\subsection{排版与发布交接}\label{ux6392ux7248ux4e0eux53d1ux5e03ux4ea4ux63a5}

排版参照用户指定的《利用初等微分方程与围道积分的傅里叶变换反演(1).pdf》，并沿用前三章的
LaTeX 模板：A4、11pt、23mm 页边距、1.12 行距；中文 Fandol
宋体、标题加粗宋体、数学 Computer Modern/AMS。目录、47
条编号索引及逐节索引均设置内部跳转，习题另起部分。

正文按现有网站结构整理为五节 Markdown
与独立习题页，附件路径及相邻节导航齐全。最终发布时仍须在实际仓库中运行既定检查、测试、构建、链接验收与部署核验；本次没有把本地
PDF 校审代替网站发布验收。



\subsection{2026 年 11 月 20 日后续修订}
推进第五章时补全原函数级数收敛半径的反证，Abel 定理改用显式包含端点的紧扇形，习题 4.2.10 的面积求和按教材从 \(n=1\) 开始；另补 Weierstrass 判别法逆命题的反例。上述改动已同步至网站稿；本段保留原有校审记录的历史含义，不把本次局部修订视为全章重新逐项验收。


\end{document}
